\documentclass[12pt, a4paper, onecolumn, twoside]{report}
% File laporan_setting.tex harus Anda sesuaikan dengan informasi Anda.
%
% Template skripsi / tesis ini berdasarkan aturan yang dikeluarkan UI di tahun 2010.
% Yang mengembangkan template ini dari awal hingga saat ini:
% 1. Fahrurrozi Rahman.
% 2. Adreas Febrian, Lia Sadita, Andre Tampubolon, Erik Dominikus.
% 3. Andrew, Daniel Salim (Mat UI 2010).
% 4. Mia Vania (Mat UI 2011).
% 5. Hengki Tasman
% 6. Barry (Mat UI 2010), Maulana (Mat UI 2011)
% 7. Antonio Kevin (Mat UI 2014)*

%*UPDATE: KEPUTUSAN REKTOR NOMOR 2143/SK/R/UI/2017

%NOTE! Remember to do use [\vspace*{-0.25cm}] for enumerate or other listing, so the spacing will be consistent. 
%TO DO NEXT. Change footer font into Arial.

\usepackage[paper=a4paper,headheight=30pt,left=3cm,top=3cm,right=3cm,bottom=3cm]{geometry}
\usepackage{graphicx}
\usepackage{enumerate}
\usepackage{amsmath,amssymb,amsthm}
\usepackage{mathrsfs}
\usepackage{indentfirst} % Kalimat pertama masuk beberapa cm
\usepackage{natbib} % Untuk merujuk pustaka
\usepackage{tocbibind}
%\usepackage{multirow} % Membuat banyak baris dalam 1 baris di tabel
%\usepackage{colortbl} % Membuat tabel berwarna-warni.
\usepackage{caption}
\usepackage{float}
\usepackage{times}
\usepackage{multirow}
\usepackage{subfig}
\usepackage{makecell}
\usepackage{fancyhdr} % Membuat header dan footer pada dokumen.
%\usepackage{pdfpages} % Package untuk memasukan berkas pdf sebagai bagian dari dokumen.

%\usepackage{colortbl} % Membuat tabel berwarna-warni.
\usepackage{algorithmicx}
\usepackage{algcompatible}

\usepackage{hyperref}
\hypersetup{a4paper,pdfstartview=FitH,pdfview=FitH}

\usepackage{array}
\usepackage{url}
\urlstyle{same}

\usepackage{setspace}
\linespread{1.3} % Jarak baris 1,5

\usepackage[ConnyRevised]{fncychap} % Membantu dalam mengatur jarak antara tepi kertas dengan posisi header.


\usepackage{threeparttable}

%-----------------------------------------------------------------------------%
% Konfigurasi
%-----------------------------------------------------------------------------%
\sloppy

\newif\ifpdf
\ifx\pdfoutput\undefined
   \pdffalse
\else
   \pdfoutput=1
   \pdftrue
\fi

% link page numbers in TOC
\makeatletter
\def\contentsline#1#2#3#4{%
  \ifx\\#4\\%
    \csname l@#1\endcsname{#2}{#3}%
  \else
    \csname l@#1\endcsname{%
      \hyper@linkstart{link}{#4}{#2}\hyper@linkend
    }{%
      % same link destination for the page:
      \hyper@linkstart{link}{#4}{#3}\hyper@linkend
      % link destination is the page itself:
      % \hyperpage{#3}%
    }%
  \fi
}

\renewcommand\section{\@startsection {section}{1}{\z@}%
                                  {-3.5ex \@plus -1ex \@minus -.2ex}%
                                  {2.3ex \@plus.2ex}%
                                  {\normalsize \bfseries}}

\renewcommand\subsection{\@startsection{subsection}{2}{\z@}%
                                     {-3.5ex\@plus -1ex \@minus -.2ex}%
                                     {1.5ex \@plus .2ex}%
                                     {\normalsize\bfseries}}
\makeatother

\setcounter{secnumdepth}{3}

\setcounter{tocdepth}{3}

\newcommand{\f}[1]{\textit{#1}}
\newcommand{\bi}[1]{\textbf{\textit{#1}}}
\newcommand{\bo}[1]{\textbf{#1}}

% Perintah untuk membuat perintah/variabel baru.
\newcommand{\var}[2]{\newcommand{#1}{#2}}

% Perintah untuk membuat perintah/variabel baru. Teks yang ditulis dalam
% perintah ini akan diformat ulang menggunakan huruf kapital.
\newcommand{\Var}[2]{\newcommand{#1}{\uppercase{#2}}}

% Tambahkan kata-kata yang dimasukan kedalam Table of Contents.
\newcommand{\addChapter}[1]{\phantomsection \addcontentsline{toc}{chapter}{#1}}

\renewcommand{\bibname}{\normalsize DAFTAR REFERENSI}
\renewcommand{\contentsname}{\normalsize DAFTAR ISI}
\renewcommand{\listfigurename}{\normalsize DAFTAR GAMBAR}
\renewcommand{\listtablename}{\normalsize DAFTAR TABEL}
\renewcommand{\figurename}{\normalsize Gambar}
\renewcommand{\tablename}{\normalsize Tabel}
\renewcommand{\chaptername}{\normalsize BAB}
\def\proofname{Bukti}

\theoremstyle{definition} % Isi teorema tidak ditulis miring
\newtheorem{theorem}{Teorema}[chapter]
\newtheorem{lemma}[theorem]{Lema} 
\newtheorem{proposition}[theorem]{Proposisi}
\newtheorem{corollary}[theorem]{Akibat}
\newtheorem{definition}{Definisi}[chapter]
\newtheorem{example}{Contoh}[chapter]

\newenvironment{remark}[1][Catatan]{\begin{trivlist}
\item[\hskip \labelsep {\normalsize \bfseries #1}]}{\end{trivlist}}

\usepackage{titlesec}
\titlespacing*{\chapter}{0pt}{-32pt}{20pt}
\titleformat{\chapter}[display]{\center \normalsize \bfseries}{\chaptertitlename \, \thechapter}{0pt}{\normalsize} % Memperbaiki judul tiap bab

% Memuat konfigurasi khusus untuk laporan yang sedang dibuat
%-----------------------------------------------------------------------------
% Informasi Mengenai Dokumen
%-----------------------------------------------------------------------------

% Judul laporan.
% Tulis judul Skripsi/Thesis/Disertasi Anda di sini
\var{\judul}{Analisis Kestabilan Global pada Model Infeksi HIV}

% Tulis kembali judul laporan, kali ini akan diubah menjadi huruf kapital secara otomatis
% Tulis judul Skripsi/Thesis/Disertasi Anda di sini
\Var{\Judul}{Analisis Kestabilan Global pada Model Infeksi HIV}

% Tulis kembali judul laporan namun dengan bahasa Ingris
\var{\judulInggris}{Global Stability Analysis of an HIV Infection Model}

% Tipe laporan, dapat berisi Skripsi, Tesis atau Disertasi
\var{\type}{Skripsi / Tesis / Disertasi} % Pilih satu, hapus yang tidak perlu
\Var{\Type}{Skripsi / Tesis / Disertasi} % Pilih satu, hapus yang tidak perlu, Akan jadi huruf besar semua secara otomatis

% Tulis nama penulis
\var{\penulis}{Nurdini K.}
\Var{\Penulis}{Nurdini K.} % Akan jadi huruf besar semua secara otomatis

% Tulis NPM penulis
\var{\npm}{1006673613}

% Tuliskan Fakultas dimana penulis berada
\Var{\Fakultas}{Matematika dan Ilmu Pengetahuan Alam}
\var{\fakultas}{Matematika dan Ilmu Pengetahuan Alam}

% Tuliskan Program Studi yang diambil penulis
\Var{\Program}{Sarjana Matematika / Statistika / Ilmu Aktuaria}
\var{\program}{Sarjana Matematika / Statistika / Ilmu Aktuaria}

% Tuliskan Program Studi yang diambil penulis, dalam bahasa Inggris
\Var{\Programinggris}{Undergraduate Study Program of Mathematics / Statistics / Actuarial Science}
\var{\programinggris}{Undergraduate Study Program of Mathematics / Statistics / Actuarial Science}

% Tuliskan bulan dan tahun publikasi laporan
\Var{\bulanTahun}{Juni 2025}
\var{\tahun}{2025}

% Tuliskan gelar yang akan diperoleh dengan menyerahkan laporan ini
\var{\gelar}{Sarjana / Magister / Doktor} % Pilih satu, hapus yang tidak perlu

% Tuliskan tanggal pengesahan laporan, waktu dimana laporan diserahkan ke
% penguji/sekretariat
\var{\tanggalPengesahan}{17 Juli 2025}

% Tuliskan tanggal keputusan sidang dikeluarkan dan penulis dinyatakan lulus/tidak lulus
\var{\tanggalLulus}{17 Juli 2025}

% Tuliskan nama pembimbing
\var{\pembimbinga}{Dr. Hengki Tasman}
\var{\pembimbingb}{Nama Pembimbing II} % Beri tanda % di awal baris ini jika Anda tidak punya Pembimbing II skripsi

%Tuliskan nama penguji
\var{\pengujia}{Nama Penguji I}
\var{\pengujib}{Nama Penguji II}
 % Anda HARUS edit file laporan_setting.tex

% Memuat Daftar pemenggalan suku kata dan istilah dalam LaTeX
%
% Hyphenation untuk Indonesia
%
% @author  Andreas Febrian
% @version 1.00
%
% Tambahkan cara pemenggalan kata-kata yang salah dipenggal secara otomatis
% oleh LaTeX. Jika kata tersebut dapat dipenggal dengan benar, maka tidak
% perlu ditambahkan dalam berkas ini. Tanda pemenggalan kata menggunakan
% tanda '-'; contoh:
% menarik
%   --> pemenggalan: me-na-rik
%

\hyphenation{
    % alphabhet A
    aki-bat
    ana-li-sis
    ana-li-tik
    ang-go-ta
    an-ti-bio-tik
    % alphabhet B
    ba-gai-ma-na
    be-be-ra-pa
    ber-ge-rak
    ber-hu-bu-ngan
    ber-pe-nga-ruh
    ber-te-ri-ma
    bi-la-ngan
    bim-bi-ngan
    % alphabhet C
    % alphabhet D
    de-fi-ni-si
    de-ngan
    di-ban-ding-kan
    dia-gram
    di-ba-ngun
    di-buk-ti-kan
    di-da-pat-kan
    di-de-fi-ni-si-kan
    di-fe-ren-si-al
    di-ha-rap-kan
    di-li-hat
    di-mo-del-kan
    di-mu-lai
    di-nya-ta-kan
    di-pe-nga-ru-hi
    di-sa-ji-kan
    di-re-pre-sen-ta-si-kan
    di-ting-kat-kan
    du-ku-ngan
    % alphabhet E
    eks-pli-sit
    eks-pre-si
    eva-lua-si
    % alphabhet F
    frak-sio-nal
    % alphabhet G
    geo-me-tri
    % alphabhet H
    ha-rap
    hi-tung
    % alphabhet I
    im-pli-sit
    in-te-gral
    ite-ra-si
    ite-ra-tif
    % alphabhet J
    % alphabhet K
    ka-it
    ke-mung-kin-an
    ke-se-im-ba-ngan
    ke-se-tim-ba-ngan
    ki-rim
    kom-par-te-men
    kon-fi-gu-ra-si
    kon-struk-si
    % alphabhet L
    li-mit
    li-te-ra-tur
    % alphabhet M
    ma-te-ma-ti-ka
    me-mo-del-kan
    me-mo-ri
    men-de-ngar-kan
    me-ne-ri-ma
    me-nga-it-kan
    me-nga-ju-kan
    me-nga-la-mi
    me-nga-na-li-sis
    me-nga-pa
    me-nga-rah-kan
    meng-ar-ti-kan
    me-nge-nai
    me-ngi-rim-kan
    me-nya-ta-kan
    mi-sal
    mi-sal-kan
    mo-del
    mu-lai
    mung-kin
    % alphabhet N
    na-pas
    non-eks-klu-sif
    % alphabhet O
    ope-ra-tor
    % alphabhet P
    pa-ling
    pa-ra-me-ter
    pa-sang
    pa-sang-an
    pe-mo-del-an
    pe-nga-ruh
    pe-nga-wa-san
    pe-ngem-ba-ngan
    pe-ngo-bat-an
    pen-ting
    pe-ngu-ra-ngan
    pe-nye-le-sai-an
    per-ge-rak-an
    per-hi-tung-an
    pe-ri-la-ku
    per-ten-ta-ngan
    per-u-mum-an
    peru-ba-han
    po-pu-la-si
    pro-duk-ti-vi-tas
	% alphabhet Q
    % alphabhet R
    re-pre-sen-ta-si
    % alphabhet S
    sa-ngat
    sa-tu
    sa-tu-nya
    se-dang
    se-dang-kan
    se-ki-tar
    se-ring
    si-mu-la-si
    su-a-tu
    sub-sti-tu-si
    % alphabhet T
    ten-tang
    te-ra-pi
    te-ri-ma
    tri-li-un
    tu-run-an
    % alphabhet U
    % alphabhet V
    % alphabhet W
    % alphabhet X
    % alphabhet Y
    % alphabhet Z
    % special
}  % Tambahkan daftar pemenggalan kata dalam file hype.indonesia.tex

% Memuat Daftar istilah yang mungkin perlu ditandai
%
% @author  Andreas Febrian
% @version 1.00
% 
% Mendaftar seluruh istilah yang mungkin akan perlu dijadikan 
% italic atau bold pada setiap kemunculannya dalam dokumen. 
% 

\var{\license}{\f{Creative Common License 1.0 Generic}}
\var{\bslash}{$\setminus$}

\begin{document}

% Sampul Depan Skripsi / Tesis
%
% Sampul Laporan

%
% @author  unknown
% @version 1.01
% @edit by Andreas Febrian
%

\begin{titlepage}
    \begin{center}    
        \begin{figure}
            \begin{center}
                \includegraphics[width=2.5cm]{pics/makara.png}
            \end{center}
        \end{figure}    
        \vspace*{0cm}
        \bo{
        	UNIVERSITAS INDONESIA\\
        }
        
        \vspace*{1.0cm}
        % judul thesis harus dalam 14pt Times New Roman
        \bo{\Judul} \\[1.0cm]

        \vspace*{2.5 cm}    
        % harus dalam 14pt Times New Roman
        \bo{\Type}

        \vspace*{3 cm}       
        % penulis dan npm
        \bo{\Penulis} \\
        \bo{\npm} \\

        \vspace*{5.0cm}

        % informasi mengenai fakultas dan program studi
        \bo{
        	FAKULTAS \Fakultas\\
        	PROGRAM STUDI \Program \\
        	DEPOK \\
        	\bulanTahun
        }
    \end{center}
\end{titlepage}
 % Dilarang mengubah isi file sampul.tex 

\pagenumbering{roman} % Gunakan penomoran halaman romawi

%pengaturan headerfooterhalaman
\renewcommand{\headrulewidth}{0.0pt}
	\fancyhf{}
	\fancyhead[L]{}
	\fancyhead[C]{}
	\fancyhead[R]{}
	\renewcommand{\headrulewidth}{0.0pt}
	\fancyfoot[C]{\thepage}
	\renewcommand{\footrulewidth}{0pt}
\pagestyle{fancy}

% Memuat halaman judul dalam
\addChapter{HALAMAN JUDUL}
%
% Halaman Judul Laporan 
%
% @author  unknown
% @version 1.01
% @edit by Andreas Febrian
%

\begin{titlepage}
    \begin{center}\begin{figure}
            \begin{center}
                \includegraphics[width=2.5cm]{pics/makara.png}
            \end{center}
        \end{figure}    
        \vspace*{0cm}
        \bo{
        	UNIVERSITAS INDONESIA\\
        }
        
        \vspace*{1.0cm}
        % judul thesis harus dalam 14pt Times New Roman
        \bo{\Judul} \\[1.0cm]

        \vspace*{2.5 cm}    
        % harus dalam 14pt Times New Roman
        \bo{\Type} \\
        % keterangan prasyarat
        \bo{Diajukan sebagai salah satu syarat untuk memperoleh gelar \\
        \gelar}\\

        \vspace*{3 cm}       
        % penulis dan npm
        \bo{\Penulis} \\
        \bo{\npm} \\

        \vspace*{3.5cm}

        % informasi mengenai fakultas dan program studi
        \bo{
        	FAKULTAS \Fakultas\\
        	PROGRAM STUDI \Program \\
        	DEPOK \\
        	\bulanTahun
        }
    \end{center}
\end{titlepage} % Dilarang mengubah isi file judul_dalam.tex

\setcounter{page}{2} % nomor halaman selanjutnya adalah 2

% Memuat halaman pengesahan
%\addChapter{LEMBAR PERSETUJUAN}
%%
% Halaman Pengesahan
%
% @author  Andreas Febrian
% @version 1.01
%

\chapter*{\normalsize HALAMAN PERSETUJUAN}

\vspace*{0.2cm}
\noindent

\noindent
\begin{tabular}{l l p{11cm}}
	Judul&: & \judul \\
	Nama&: & \penulis \\
	NPM&: & \npm \\
\end{tabular} \\

\vspace*{1.2cm}

\noindent Laporan \type~ini telah diperiksa dan disetujui.\\[0.3cm]
\begin{center}
\tanggalPengesahan \\[2cm]


\underline{\pembimbinga}\\[0.1cm]
Pembimbing \type
\end{center}

\newpage  % Dilarang mengubah isi file pengesahan.tex 

% load halaman orisinalitas
\addChapter{HALAMAN PERNYATAAN ORISINALITAS}
%
% Halaman Orisinalitas
%
% @author  Andreas Febrian
% @version 1.01
%

\chapter*{\normalsize HALAMAN PERNYATAAN ORISINALITAS}
\vspace*{2cm}
\begin{doublespace}
\begin{center}
	\type~ini adalah hasil karya saya sendiri, \\
	dan semua sumber baik yang dikutip maupun dirujuk \\
	telah saya nyatakan dengan benar. \\
	\vspace*{2.6cm}
	
	\begin{tabular}{l c l}
	Nama & : & \penulis \\
	NPM & : & \npm \\
	Tanda Tangan & : & \\
	& & \\
	& & \\
	Tanggal & : & \tanggalPengesahan \\	
	\end{tabular}
\end{center}
\end{doublespace}

\newpage  % Dilarang mengubah isi file orisinal.tex

%pengaturan rata kiri
%\raggedright
\setlength\parindent{0.7cm}

\addChapter{HALAMAN PENGESAHAN}
%
% Halaman Pengesahan Sidang
%
% @author  Andreas Febrian, Andre Tampubolon
% @version 1.02
%

\chapter*{\normalsize HALAMAN PENGESAHAN}
\begin{singlespace}
\vspace*{0.4cm}
\noindent

\noindent
\begin{tabular}{ll p{9cm}}
	\type~ini diajukan oleh&: & \\ % ---- Jangan diedit ----
	Nama&: & \penulis \\ % ---- Jangan diedit ----
	NPM&: & \npm \\ % ---- Jangan diedit ----
	Program Studi&: & \program \\ % ---- Jangan diedit ----
	Judul \type&: & \judul \\ % ---- Jangan diedit ----
\end{tabular} \\

\vspace*{1.0cm}

\noindent Telah berhasil dipertahankan di hadapan Dewan Penguji
dan diterima sebagai bagian persyaratan yang diperlukan untuk
memperoleh gelar \gelar \, pada Program Studi \program, Fakultas
\fakultas, Universitas Indonesia.\\[0.2cm] % ---- Jangan diedit ----

\begin{center}
	DEWAN PENGUJI
\end{center}

\vspace*{0.3cm}

\begin{tabular}{l l l l }
& & & \\
Pembimbing I &: & \pembimbinga & (\hspace*{3.0cm}) \\ % ---- Jangan diubah ----

\vspace{0.7cm}

& & & \\


Pembimbing II&: & \pembimbingb & (\hspace*{3.0cm}) \\ % ---- Silakan hapus bagian ini jika tidak ada pembimbing II

\vspace{0.7cm}

& & & \\


Penguji I&: & \pengujia & (\hspace*{3.0cm}) \\ % ---- Jangan diubah ------

\vspace{0.7cm}

& & & \\


Penguji II&: & \pengujib & (\hspace*{3.0cm}) \\ % ---- Jangan diubah ------
& & & \\
\end{tabular}\\

\vspace*{2.0cm}

\begin{tabular}{ll l}
	Ditetapkan di&: & Depok\\ % ---- Jangan diedit ----
	Tanggal&: & \tanggalLulus \\ % ---- Jangan diedit ----
\end{tabular}
\end{singlespace}

\newpage  % File pengesahan_sidang.tex mungkin perlu diubah sesuai dengan banyaknya pembimbing dan penguji kolokium atau sidang Anda

\addChapter{KATA PENGANTAR}
%-----------------------------------------------------------------------------%
\chapter*{\normalsize KATA PENGANTAR}
%-----------------------------------------------------------------------------%

Puji syukur kepada Tuhan Yang Maha Esa atas segala berkat dan rahmat-Nya, sehingga penulis dapat menyelesaikan skripsi ini dengan sebaik-baiknya. Selama masa penulisan skripsi ini, penulis telah mendapat banyak dukungan, doa, bantuan, inspirasi, dan motivasi dari berbagai pihak. Untuk itu, penulis ingin berterima kasih kepada:
\begin{enumerate}
	\item Semua staf pengajar atas ilmu pengetahuan yang telah diberikan kepada penulis selama masa kuliah.
	\item Seluruh staf karyawan yang selalu melakukan tugas mereka dengan baik, sehingga memberikan kenyamanan pelayanan bagi siapa saja.
	\item Seluruh teman-teman angkatan 2009 atas semangat, canda tawa dan banyak hal penting selama masa kuliah.
	\item Kakak-kakak angkatan 2008, 2007, dan 2006 serta adik-adik angkatan 2010, dan 2011 yang terus memberi semangat.
\end{enumerate}

\vspace*{0.1cm}

\begin{flushright}
\penulis \\[0.1cm]
\vspace*{0.2cm}
\tahun
\end{flushright}  % Anda HARUS isi file pengantar.tex -------------

\addChapter{HALAMAN PERSETUJUAN PUBLIKASI ILMIAH}
%
% @author  Andre Tampubolon, Andreas Febrian
% @version 1.01
%

\chapter*{\normalsize HALAMAN PERNYATAAN PERSETUJUAN PUBLIKASI TUGAS AKHIR UNTUK KEPENTINGAN AKADEMIS}
\begin{singlespace}
\vspace*{0.2cm}
\noindent
Sebagai sivitas akademik Universitas Indonesia, saya yang bertanda
tangan di bawah ini:
\vspace*{0.4cm}


\begin{tabular}{p{3cm} l p{8cm}}
	Nama & : & \penulis \\ 	
	NPM & : & \npm \\
	Program Studi & : & \program\\	
	Fakultas & : & \fakultas\\
	Jenis Karya & : & \type \\
\end{tabular}

\vspace*{0.6cm}
\noindent demi pengembangan ilmu pengetahuan, menyetujui untuk memberikan
kepada Universitas Indonesia Hak Bebas Royalti Noneksklusif
(\textit{Non-exclusive Royalty Free Right}) atas karya ilmiah saya yang berjudul:
\begin{center}
	\judul
\end{center}
beserta perangkat yang ada (jika diperlukan). Dengan Hak Bebas Royalti
Noneksklusif ini, Universitas Indonesia berhak menyimpan,
mengalihmedia/formatkan, mengelola dalam bentuk pangkalan data
(\textit{data base}), merawat, dan memublikasikan tugas akhir saya selama
tetap mencantumkan nama saya sebagai penulis/pencipta dan sebagai
pemilik Hak Cipta. \\

\noindent Demikian pernyataan ini saya buat dengan sebenarnya.

\begin{center}
	\vspace*{0.8cm}
	\begin{tabular}{lll}
		Dibuat di&: & Depok \\
		Pada tanggal&: & \tanggalPengesahan \\
	\end{tabular}\\

	\vspace*{0.2cm}
	Yang menyatakan \\
	\vspace*{1.1cm}
	(\penulis)
\end{center}
\end{singlespace}

\newpage

 % Dilarang ubah file persetujuan_publikasi.tex 

%pengaturan header footer halaman
\renewcommand{\headrulewidth}{0.0pt}
	\fancyhf{}
	\fancyhead[L]{}
	\fancyhead[C]{}
	\fancyhead[R]{}
	\renewcommand{\headrulewidth}{0.0pt}
	\fancyfoot[C]{\thepage}
	\fancyfoot[R]{{\fontfamily{phv}\selectfont\footnotesize \bo{Universitas Indonesia}}}
	\renewcommand{\footrulewidth}{0.0pt}
\pagestyle{fancyplain}

\addChapter{ABSTRAK}
%
% Halaman Abstrak
%
% @author  Andreas Febrian
% @version 1.00
%

\chapter*{\normalsize ABSTRAK}
\begin{singlespace}
\vspace*{0.2cm}

\noindent\begin{tabular}{l l p{10cm}}
Nama&: & \penulis \\ % Jangan diubah
Program Studi&: & \program \\ % Jangan diubah
Judul&: & \judul \\ % Jangan diubah
\end{tabular}

\vspace*{0.5cm}

\noindent Secara matematis, melipat dapat dilakukan dengan merotasi bidang kertas yang ingin dilipat terhadap sumbu garis lipatan. Pemetaan dari kertas ke hasil lipatan origami dilakukan dengan merotasi bidang kertas yang sesuai. Apabila sebuah origami yang telah selesai dibuka kembali, terdapat garis-garis bekas lipatan pada kertas. Garis-garis lipatan ini disebut sebagai pola lipatan. Jika sebuah pola lipatan dapat dilipat, perkalian matriks-matriks rotasi sesuai merupakan matriks identitas. Hal ini berlaku pada origami simpul tunggal, dan berlaku secara lokal pada origami simpul jamak, namun dapat diperluas sehingga berlaku secara global pada origami simpul jamak.\\

\vspace*{0.2cm}

\noindent Kata kunci:\\
	Origami, rotasi, transformasi, lipatan tak datar. \\ % ------------ Jangan lupa diedit
\end{singlespace}
\newpage  % Abstrak dalam bahasa Indonesia. Ubah file abstrak.tex.

\addChapter{ABSTRACT}
%
% Halaman Abstract
%
% @author  Andreas Febrian
% @version 1.00
%

\chapter*{\normalsize ABSTRACT}
\begin{singlespace}
\vspace*{0.2cm}

\noindent \begin{tabular}{l l p{11.0cm}}
	Name&: & \penulis \\ % ------------ Jangan diedit
	Program&: & \programinggris \\ % ------------ Jangan diedit
	Title&: & \judulInggris \\ % ------------ Jangan diedit
\end{tabular} \\

\vspace*{0.5cm}

\noindent Mathematically, to fold a paper is to rotate the paper along a crease line as the axis. The mapping from the paper to the finished origami fold is done by rotating parts of the paper to the appropriate locations. Unfolding finished origami reveals a pattern of crease lines, known as crease pattern. If the crease pattern is foldable, then the product 
of the associated rotational matrices is the identity matrix. This condition holds in a single vertex crease pattern and holds locally in a multiple vertex crease pattern and can be adapted to a global condition in a multiple vertex crease pattern\\

\vspace*{0.2cm}

\noindent Keywords:\\
	Origami, rotation, transformation, non-flat folding.\\ % ------------ Jangan lupa diedit
\end{singlespace}

\newpage  % Abstrak dalam bahasa Inggris. Ubah file abstract.tex.

% Daftar isi, gambar, dan tabel
\tableofcontents
\clearpage

\listoftables % jika ada tabel
\clearpage

\listoffigures % jika ada gambar
\clearpage{\thispagestyle{empty}\cleardoublepage}

% Gunakan penomoran Arab (1, 2, 3, ...) setelah bagian ini.
\pagenumbering{arabic}
%pengaturan headerfooter halaman
\renewcommand{\headrulewidth}{0.0pt}
   \fancyhf{}
	\fancyhead[RO,LE]{\thepage}
	\renewcommand{\headrulewidth}{0.0pt}
	\fancyfoot[R]{\footnotesize \bo{Universitas Indonesia}}
	\renewcommand{\footrulewidth}{0.0pt}
\pagestyle{fancy}

\fancypagestyle{plain}{ %
  \fancyhf{} 
  	\fancyhead[L]{}
  	\fancyhead[C]{}
  	\fancyhead[R]{}
  	\renewcommand{\headrulewidth}{0pt} 
  	\fancyfoot[C]{\thepage}
  	\fancyfoot[R]{\footnotesize \bo{Universitas Indonesia}}
  	\renewcommand{\footrulewidth}{0pt}
}

\clearpage{\thispagestyle{empty}\cleardoublepage}
\chapter{Pendahuluan}

\section{Latar Belakang}

Peningkatan harapan hidup global telah menjadi salah satu pencapaian penting dalam bidang kesehatan dan sosial pada abad ke-21. Fenomena ini ditandai dengan bertambahnya jumlah individu yang mencapai usia sangat lanjut, termasuk kelompok supercentenarian yaitu individu berusia 110 tahun atau lebih \citep{Young2021}. Data menunjukkan bahwa jumlah supercentenarian yang tervalidasi meningkat secara signifikan, dengan kasus pertama tercatat pada tahun 1960 dan terus berkembang hingga saat ini. Kondisi ini membawa implikasi penting bagi sistem jaminan sosial, industri asuransi jiwa, dan dana pensiun yang harus mengantisipasi risiko umur panjang (\textit{longevity risk}) dengan lebih cermat.

Risiko umur panjang mengacu pada ketidakpastian terkait dengan proyeksi harapan hidup yang lebih tinggi dari perkiraan semula. Ketika individu hidup lebih lama dari yang diperkirakan, institusi keuangan yang memberikan jaminan pembayaran jangka panjang seperti anuitas dan dana pensiun menghadapi beban finansial yang lebih besar. Oleh karena itu, pemodelan mortalitas pada usia lanjut yang akurat menjadi kebutuhan mendasar untuk mengukur dan mengelola risiko ini secara efektif.

Model mortalitas klasik seperti Gompertz-Makeham, Heligman-Pollard, dan model logistik telah lama digunakan dalam praktik aktuaria untuk memodelkan pola kematian manusia. Namun, model-model tersebut menunjukkan keterbatasan dalam menangkap fenomena deselerasi mortalitas pada usia lanjut ekstrem, yaitu kondisi di mana laju peningkatan tingkat kematian melambat atau bahkan mendekati konstan pada usia sangat tinggi \citep{Thatcher1999}. Selain itu, model-model klasik umumnya menetapkan batas usia maksimum tabel mortalitas (\textit{highest attained age}, $\omega$) secara subjektif tanpa landasan statistik yang kuat, sehingga berpotensi menghasilkan estimasi yang kurang reliabel pada ekor distribusi usia kematian.

Untuk mengatasi keterbatasan tersebut, pendekatan berbasis \textit{extreme value theory} (EVT) mulai diterapkan dalam pemodelan mortalitas usia lanjut. EVT merupakan cabang statistik yang secara khusus dirancang untuk menganalisis perilaku ekor distribusi dan kejadian ekstrem. Dalam konteks mortalitas, EVT memberikan kerangka kerja yang lebih sesuai untuk memodelkan individu yang mencapai usia sangat lanjut dengan menggunakan \textit{generalized Pareto distribution} (GPD). Salah satu model yang mengadopsi EVT adalah \textit{threshold life table} (TLT), yang dikembangkan untuk menggabungkan model Gompertz pada usia non-ekstrem dengan GPD pada usia lanjut ekstrem, dengan pemisahan kedua komponen dilakukan pada suatu usia ambang (\textit{threshold age}, $N$).

Meskipun model TLT memberikan fleksibilitas dalam memodelkan mortalitas usia lanjut, model ini memiliki kelemahan potensial berupa ketidakmulusingan (\textit{discontinuity}) pada fungsi \textit{hazard} di titik usia ambang. Ketidakmulusingan ini dapat menimbulkan interpretasi yang kurang realistis secara biologis dan mengurangi kualitas proyeksi mortalitas. Untuk mengatasi masalah tersebut, dikembangkan model \textit{smoothed threshold life table} (STLT) yang menambahkan kendala kehalusan (\textit{smoothing constraint}) sehingga memastikan transisi yang mulus antara komponen Gompertz dan GPD pada usia ambang.

Lebih lanjut, untuk mengakomodasi perubahan pola mortalitas antar kohor dan memungkinkan peramalan mortalitas ke masa depan, model STLT diperluas menjadi model dinamis yang disebut \textit{dynamic smooth threshold life table} (DSTLT). Model ini memperkenalkan komponen waktu dengan memodelkan parameter tertentu sebagai fungsi dari indeks kohor, sehingga mampu menangkap tren mortalitas jangka panjang dan melakukan proyeksi untuk kohor-kohor mendatang.

Penelitian ini bertujuan untuk mengkaji secara mendalam konstruksi, estimasi parameter, dan evaluasi kinerja model STLT dan DSTLT dalam konteks pemodelan mortalitas usia lanjut. Kajian ini diharapkan dapat memberikan kontribusi pada pemahaman yang lebih baik mengenai penerapan EVT dalam aktuaria, serta menyediakan alternatif model yang lebih robust untuk mengukur dan mengelola risiko umur panjang dalam industri asuransi dan dana pensiun.

\section{Rumusan Masalah}

Berdasarkan latar belakang yang telah diuraikan, penelitian ini difokuskan pada kajian mendalam terhadap model STLT dan DSTLT dalam konteks pemodelan mortalitas usia lanjut. Rumusan masalah penelitian ini dijabarkan dalam pertanyaan-pertanyaan berikut:

\begin{enumerate}
    \item Bagaimana konstruksi matematis model STLT dan DSTLT, khususnya terkait penambahan kendala kehalusan pada fungsi \textit{hazard} dan pemodelan parameter dinamis?
    
    \item Bagaimana prosedur estimasi parameter untuk model STLT dan DSTLT menggunakan metode \textit{maximum likelihood estimation} (MLE) dengan mempertimbangkan karakteristik data mortalitas yang tersensor interval dan kanan, serta penentuan usia ambang optimal?
    
    \item Bagaimana kinerja model STLT dibandingkan dengan model mortalitas statis lain seperti Gompertz-Makeham, Heligman-Pollard, dan Coale-Kisker dalam hal kesesuaian terhadap data historis?
    
    \item Bagaimana kinerja model DSTLT dalam hal kesesuaian \textit{in-sample} dan kemampuan peramalan \textit{out-of-sample}, khususnya dibandingkan dengan model mortalitas dinamis yang relevan seperti model Cairns-Blake-Dowd (CBD)?
\end{enumerate}

\section{Tujuan Penelitian}

Berdasarkan rumusan masalah yang telah dipaparkan, penelitian ini memiliki tujuan-tujuan sebagai berikut:

\begin{enumerate}
    \item Mengonstruksi model STLT dan DSTLT secara matematis, termasuk derivasi kendala kehalusan pada fungsi \textit{hazard} untuk model STLT dan formulasi parameter dinamis untuk model DSTLT.
    
    \item Melakukan estimasi parameter model STLT dan DSTLT menggunakan metode \textit{maximum likelihood estimation} (MLE) dengan mempertimbangkan data mortalitas tersensor interval dan kanan, serta menentukan usia ambang optimal.
    
    \item Mengevaluasi kinerja model STLT dengan membandingkan kesesuaian model (\textit{goodness-of-fit}) terhadap model mortalitas statis lain seperti Gompertz-Makeham, Heligman-Pollard, dan Coale-Kisker.
    
    \item Mengevaluasi kinerja model DSTLT dalam hal kesesuaian \textit{in-sample} dan kemampuan peramalan \textit{out-of-sample}, dengan membandingkannya terhadap model mortalitas dinamis Cairns-Blake-Dowd (CBD).
\end{enumerate}

\section{Batasan Penelitian}

 Penelitian ini mengasumsikan tidak terjadi migrasi pada populasi yang diteliti. Asumsi ini diperlukan untuk menyederhanakan analisis kohor dan memastikan perubahan jumlah individu dalam kohor hanya disebabkan oleh kematian. Untuk memenuhi asumsi ini, dilakukan penyesuaian data untuk menciptakan kohor tertutup yang menghilangkan pengaruh migrasi.
\clearpage{\thispagestyle{empty}\cleardoublepage}
\chapter{Landasan Teori}

Bab ini berisikan teori yang menjadi dasar penelitian mengenai pemodelan mortalitas usia lanjut dengan pendekatan \textit{Dynamic Smooth Threshold Life Table} (DSTLT). Pembahasan dimulai dengan konsep dasar tabel mortalitas dan notasi aktuaria yang digunakan secara luas dalam analisis mortalitas. Selanjutnya, diuraikan berbagai model mortalitas statis seperti Gompertz-Makeham, Heligman-Pollard, dan Coale-Kisker yang menjadi perbandingan dalam penelitian ini. Pemahaman terhadap \textit{Extreme Value Theory} (EVT) dan pendekatannya dalam memodelkan kejadian ekstrem menjadi landasan penting untuk memahami konstruksi model \textit{Threshold Life Table} (TLT). Model TLT kemudian dikembangkan menjadi \textit{Smoothed Threshold Life Table} (STLT) untuk mengatasi ketidakmulsan fungsi \textit{hazard} di usia ambang. Terakhir, diperkenalkan versi dinamis dari model STLT, yaitu DSTLT, yang memungkinkan pemodelan tren mortalitas antarkohor dan peramalan mortalitas jangka panjang. Seluruh teori yang dipaparkan dalam bab ini mengacu pada kerangka kerja yang dikembangkan oleh \citet{huang2020modelling} dan literatur terkait di bidang aktuaria dan statistika ekstrem.

\section{Konsep Dasar Tabel Mortalitas}

Tabel mortalitas merupakan instrumen fundamental dalam ilmu aktuaria yang menyajikan representasi matematis dari pola kematian suatu populasi. Tabel ini menyediakan informasi kuantitatif mengenai probabilitas bertahan hidup dan kematian pada berbagai usia, yang sangat esensial untuk penilaian risiko longevitas, penetapan premi asuransi jiwa, dan perhitungan kewajiban dana pensiun \citep{dickson2020actuarial}.

\subsection{Fungsi-Fungsi Dasar Mortalitas}

Misalkan $X$ adalah variabel acak kontinu yang menyatakan usia saat kematian seseorang yang baru lahir. Fungsi distribusi kumulatif (\textit{cumulative distribution function}, CDF) dari $X$ didefinisikan sebagai
\begin{equation}
    F(x) = P(X \leq x), \quad x \geq 0,
\end{equation}
yang menyatakan probabilitas seorang individu meninggal sebelum atau pada usia $x$. Fungsi probabilitas bertahan hidup (\textit{survival function}) dinyatakan sebagai
\begin{equation}
    S(x) = P(X > x) = 1 - F(x), \quad x \geq 0,
\end{equation}
yang menunjukkan probabilitas seseorang bertahan hidup melampaui usia $x$. Fungsi kepekatan probabilitas (\textit{probability density function}, PDF) dari $X$ adalah
\begin{equation}
    f(x) = \frac{dF(x)}{dx} = -\frac{dS(x)}{dx}, \quad x \geq 0,
\end{equation}
yang menggambarkan laju kematian instantan pada usia $x$.

Fungsi \textit{hazard} atau fungsi tingkat kematian (\textit{force of mortality}), yang dinotasikan dengan $h(x)$ atau $\mu(x)$, didefinisikan sebagai
\begin{equation}
    h(x) = \lim_{\Delta x \to 0^+} \frac{P(x < X \leq x + \Delta x \mid X > x)}{\Delta x} = \frac{f(x)}{S(x)} = -\frac{d \ln S(x)}{dx}.
\end{equation}
Fungsi ini mengukur tingkat kematian sesaat pada usia $x$, dengan asumsi individu tersebut telah bertahan hingga usia tersebut. Hubungan antara fungsi \textit{hazard} dan fungsi \textit{survival} dapat dinyatakan melalui relasi integral
\begin{equation}
    S(x) = \exp\left(-\int_0^x h(t) \, dt\right).
\end{equation}

\subsection{Notasi Aktuaria}

Dalam praktik aktuaria, digunakan notasi diskret yang mengadaptasi fungsi-fungsi kontinu di atas ke dalam interval usia yang terpisah, biasanya tahunan. Beberapa notasi aktuaria standar yang digunakan dalam penelitian ini adalah sebagai berikut \citep{dickson2020actuarial}:

\begin{itemize}
    \item $l_x$: Jumlah individu yang bertahan hidup hingga usia tepat $x$ dalam suatu kohor hipotetis.
    \item $d_x = l_x - l_{x+1}$: Jumlah kematian yang terjadi antara usia $x$ dan $x+1$.
    \item $q_x = \frac{d_x}{l_x}$: Probabilitas seorang individu berusia $x$ meninggal sebelum mencapai usia $x+1$.
    \item $p_x = 1 - q_x = \frac{l_{x+1}}{l_x}$: Probabilitas seorang individu berusia $x$ bertahan hidup hingga usia $x+1$.
    \item $e_x = \sum_{k=1}^{\omega - x} k \cdot {}_{k|} q_x$: Harapan hidup lengkap pada usia $x$, yaitu ekspektasi jumlah tahun yang akan dijalani oleh seseorang yang kini berusia $x$.
    \item $m_x = \frac{d_x}{L_x}$: Tingkat kematian sentral, dengan $L_x$ adalah \textit{person-years lived} antara usia $x$ dan $x+1$. Untuk interval satu tahun, sering diasumsikan $m_x \approx \frac{d_x}{(l_x + l_{x+1})/2}$.
\end{itemize}

Notasi-notasi ini memfasilitasi penghitungan dan analisis mortalitas dalam konteks aktuaria, khususnya untuk keperluan penilaian risiko dan penetapan tarif produk asuransi serta pensiun.

\subsection{Highest Attained Age dan Interval Censoring}

Dalam data mortalitas empiris, khususnya untuk usia lanjut, informasi mengenai usia kematian yang tepat seringkali tidak tersedia. Data yang tersedia biasanya dalam bentuk usia terakhir ulang tahun (\textit{last birthday age}) atau usia terakhir yang dicapai (\textit{highest attained age}). Hal ini menyebabkan data mortalitas bersifat tersensor interval (\textit{interval-censored}).

Misalkan seorang individu tercatat meninggal pada usia $k$ (usia terakhir ulang tahun). Ini berarti usia kematian sebenarnya $X$ berada dalam interval $[k, k+1)$. Dalam konteks estimasi parameter model mortalitas, fungsi likelihood harus memperhitungkan sifat tersensor ini. Untuk individu yang meninggal pada interval $[k, k+1)$, kontribusi terhadap likelihood adalah
\begin{equation}
    P(k \leq X < k+1) = S(k) - S(k+1) = S(k) \left[1 - \frac{S(k+1)}{S(k)}\right] = S(k) \cdot q_k,
\end{equation}
di mana $q_k$ adalah probabilitas kematian dalam satu tahun pada usia $k$.

Selain tersensor interval, data mortalitas usia lanjut juga sering mengalami tersensor kanan (\textit{right-censored}) untuk individu yang masih hidup pada akhir periode observasi. Jika seorang individu bertahan hidup hingga usia $k$ pada akhir pengamatan, kontribusinya terhadap likelihood adalah $S(k)$, yang mencerminkan probabilitas bertahan hingga setidaknya usia tersebut \citep{huang2020modelling}.

Pemahaman terhadap mekanisme penyensoran ini sangat penting dalam konstruksi fungsi likelihood untuk estimasi parameter model mortalitas berbasis \textit{maximum likelihood estimation} (MLE), yang akan dibahas lebih lanjut dalam bab selanjutnya.

\section{Model Mortalitas Statis}

Model mortalitas statis merupakan model parametrik yang menggambarkan pola kematian sebagai fungsi dari usia tanpa mempertimbangkan dimensi waktu atau perubahan antarkohor. Model-model ini telah digunakan secara luas dalam praktik aktuaria untuk memodelkan tabel mortalitas pada berbagai rentang usia. Bagian ini membahas beberapa model mortalitas statis yang relevan sebagai pembanding terhadap model STLT dalam penelitian ini, meliputi model Gompertz-Makeham, model Heligman-Pollard, dan metode Coale-Kisker \citep{huang2020modelling, dickson2020actuarial}.

\subsection{Model Gompertz-Makeham}

Salah satu model mortalitas parametrik paling awal dan berpengaruh dikembangkan oleh Benjamin Gompertz pada tahun 1825. Model Gompertz didasarkan pada observasi bahwa tingkat kematian dewasa cenderung meningkat secara eksponensial seiring bertambahnya usia. Fungsi \textit{hazard} atau \textit{force of mortality} dalam model Gompertz dinyatakan sebagai
\begin{equation}
    h(x) = B \exp(Cx), \quad x \geq 0,
\end{equation}
dengan $B > 0$ adalah parameter skala yang merepresentasikan tingkat mortalitas dasar, dan $C > 0$ adalah parameter bentuk yang mengatur laju peningkatan mortalitas seiring bertambahnya usia. Gompertz mengemukakan penjelasan fisiologis bahwa kapasitas seseorang untuk menghindari kematian secara bertahap menurun seiring bertambahnya usia, yang dapat dikaitkan dengan deteriorasi tubuh secara bertahap akibat akumulasi kerusakan molekuler dan seluler \citep{gompertz1825nature, thatcher1998force}.

William Makeham pada tahun 1860 menyempurnakan model Gompertz dengan menambahkan komponen konstan untuk mengakomodasi penyebab kematian yang dianggap independen terhadap usia, seperti kecelakaan atau penyakit tertentu. Fungsi \textit{hazard} dalam model Gompertz-Makeham adalah
\begin{equation}
    h(x) = A + B \exp(Cx), \quad x \geq 0,
\end{equation}
dengan $A > 0$ merepresentasikan mortalitas yang tidak terkait dengan penuaan atau pematangan biologis. Model ini dapat diinterpretasikan sebagai model kejutan (\textit{shock model}), di mana waktu hidup individu merupakan minimum dari waktu hingga kematian akibat penuaan (mengikuti distribusi Gompertz) dan waktu hingga kecelakaan fatal (mengikuti distribusi eksponensial), dengan asumsi kedua variabel acak tersebut independen \citep{makeham1860law, bower1997graduation}.

Meskipun model Gompertz-Makeham memberikan kesesuaian yang baik pada rentang usia dewasa, kelemahan utamanya adalah ketidakmampuannya menangkap fenomena deselerasi mortalitas pada usia lanjut (\textit{late-life mortality deceleration}). Fenomena ini mengacu pada perlambatan laju peningkatan mortalitas yang teramati pada usia sangat lanjut, di mana laju kematian tidak lagi meningkat secara eksponensial atau bahkan mencapai plateau \citep{barbi2018plateau, gavrilov2005mortality}. Asumsi peningkatan eksponensial tanpa batas membuat model ini kurang akurat untuk memodelkan mortalitas pada usia lanjut ekstrem, khususnya di atas 90 atau 100 tahun.

\subsection{Model Heligman-Pollard}

Model Heligman-Pollard, yang dikembangkan oleh Heligman dan Pollard pada tahun 1980, menawarkan pendekatan yang lebih komprehensif dengan menggabungkan tiga komponen utama untuk memodelkan probabilitas kematian $q_x$ pada seluruh rentang usia. Model ini berusaha menangkap pola mortalitas manusia dari kelahiran hingga usia lanjut dengan mengintegrasikan tiga sumber risiko kematian: kematian pada usia dini, kematian akibat kecelakaan, dan kematian akibat penuaan biologis. Fungsi probabilitas kematian dalam model Heligman-Pollard dinyatakan sebagai
\begin{equation}
    \frac{q_x}{1 - q_x} = A^{(x+B)^C} + D \exp\left(-E[\ln x - \ln F]^2\right) + GH^x, \quad x \geq 1,
\end{equation}
dengan $A, B, C, D, E, F, G, H$ adalah parameter model \citep{heligman1980age}. Komponen pertama $A^{(x+B)^C}$ memodelkan mortalitas pada anak usia dini yang umumnya tinggi dan menurun dengan cepat. Komponen kedua $D \exp\left(-E[\ln x - \ln F]^2\right)$ menangkap puncak mortalitas akibat kecelakaan pada usia dewasa muda. Komponen ketiga $GH^x$ merepresentasikan mortalitas akibat penuaan biologis dan dapat dipandang sebagai faktor yang berkaitan dengan hukum Gompertz.

Meskipun model Heligman-Pollard dapat memberikan kesesuaian yang baik dengan data mortalitas pada berbagai rentang usia, model ini memiliki beberapa keterbatasan penting. Pertama, komponen penuaan dalam model ini pada dasarnya mengadopsi hukum Gompertz yang, sebagaimana telah dijelaskan, tidak menggambarkan dengan baik fenomena deselerasi mortalitas pada usia lanjut ekstrem. Penelitian yang dilakukan oleh Olshansky dan Carnes pada tahun 1997 menunjukkan bahwa pola mortalitas usia lanjut tidak mengikuti hukum Gompertz \citep{olshansky1997ever}. Kedua, formulasi model menghasilkan nilai $q_x$ yang selalu kurang dari 1 untuk setiap $x > 0$. Hal ini menjadi kendala praktis dalam konstruksi tabel mortalitas lengkap, karena tabel mortalitas perlu memiliki usia maksimum di mana $q_x = 1$ atau $l_x = 0$. Dalam praktiknya, aktuaris harus menentukan usia terminal (\textit{terminal age}) secara subjektif untuk melengkapi tabel mortalitas berbasis model Heligman-Pollard \citep{huang2020modelling}.

\subsection{Metode Coale-Kisker}

Metode yang dikembangkan oleh Coale dan Kisker pada tahun 1990 merupakan metode ekstrapolasi yang sering digunakan untuk memodelkan mortalitas pada usia lanjut, khususnya di negara-negara berkembang di mana data mortalitas usia tinggi tidak tersedia atau kurang reliabel. Berbeda dengan model-model sebelumnya yang memodelkan fungsi \textit{hazard} atau probabilitas kematian secara langsung, metode Coale-Kisker bekerja dengan mengekstrapolasi tingkat kematian sentral (\textit{central death rate}), $m_x$ \citep{coale1990defects}.

Metode ini menggunakan persamaan rekursif untuk $k(x) = \ln(m_x / m_{x+1})$, yaitu logaritma natural dari rasio tingkat kematian sentral pada dua usia berurutan. Relasi rekursif tersebut adalah
\begin{equation}
    k(x) = k(x-1) - R, \quad x \geq x_0,
\end{equation}
dengan $R$ adalah konstanta penurunan dan $x_0$ adalah usia awal ekstrapolasi. Konstanta $R$ dihitung menggunakan persamaan
\begin{equation}
    R = \frac{(x_1 - x_0) k(x_0) + \ln m_{x_0} - \ln m_{x_1}}{1 + 2 + \cdots + (x_1 - x_0)},
\end{equation}
di mana $x_1$ adalah usia akhir ekstrapolasi yang dipilih. Setelah $R$ diperoleh, nilai $k(x)$ dapat dihitung secara iteratif untuk $x > x_0$, dan kemudian $m_x$ dapat dipulihkan menggunakan hubungan $m_x = m_{x-1} \exp(-k(x-1))$.

Kelemahan utama metode Coale-Kisker adalah ketergantungannya pada penetapan nilai $x_0$, $x_1$, dan $m_{x_1}$ secara subjektif. Sebagai contoh, Coale dan Kisker menggunakan $x_0 = 84$, $x_1 = 110$, dan $m_{110} = 1.0$ untuk laki-laki, serta $m_{110} = 0.8$ untuk perempuan dalam analisis mereka. Pilihan subjektif terhadap parameter-parameter ini secara langsung mempengaruhi hasil ekstrapolasi mortalitas pada usia lanjut ekstrem. Akurasi pada usia ekstrem kemungkinan rendah dengan metode ini karena kurva yang dihasilkan sangat bergantung pada titik awal dan akhir ekstrapolasi, yang merupakan kelemahan mendasar dari teknik ini \citep{huang2020modelling}.

Sebagai catatan, dalam penelitian ini, metode Coale-Kisker tidak diestimasi menggunakan \textit{maximum likelihood estimation} seperti model-model lainnya, melainkan menggunakan prosedur ekstrapolasi deterministik sebagaimana dijelaskan di atas. Oleh karena itu, perbandingan kinerja model menggunakan kriteria \textit{sum of squared errors} (SSE) antara $q_x$ observasi dan $q_x$ prediksi, bukan berdasarkan kriteria likelihood atau kriteria informasi seperti AIC atau BIC.

\section{Extreme Value Theory}

\textit{Extreme Value Theory} (EVT) merupakan cabang statistika yang secara khusus menganalisis perilaku nilai-nilai ekstrem dari suatu distribusi probabilitas. Dalam konteks mortalitas, EVT menyediakan kerangka teoretis yang kuat untuk memodelkan probabilitas kejadian langka seperti kematian pada usia lanjut ekstrem dan melakukan ekstrapolasi di luar rentang data observasi secara lebih terjustifikasi secara statistik \citep{coles2001introduction, gbari2017extreme}. Keterbatasan data empiris pada usia sangat lanjut, yang ditandai dengan jumlah observasi yang sedikit dan volatilitas tinggi, menjadikan EVT sebagai alat yang tepat karena dirancang untuk menangani kejadian ekstrem dengan fondasi teoretis yang solid.

\subsection{Konsep Dasar Extreme Value Theory}

Misalkan $T_1, T_2, \ldots, T_n$ adalah sekuens variabel acak independen yang merepresentasikan usia kematian individual dengan fungsi distribusi kumulatif yang sama $_xq_0 = P(T_i \leq x)$ untuk $x \geq 0$ dan $i = 1, \ldots, n$, dengan $_0q_0 = 0$. Variabel acak $T_i$ dapat merepresentasikan total masa hidup sejak lahir hingga kematian, atau sisa masa hidup setelah usia awal tertentu $\alpha$.

Definisikan sekuens maksimum $M_n = \max\{T_1, T_2, \ldots, T_n\}$. Dalam konteks mortalitas, $M_n$ merepresentasikan usia kematian tertinggi yang diamati dalam kelompok homogen beranggotakan $n$ individu yang tunduk pada tabel mortalitas yang sama. EVT mempelajari perilaku asimptotik dari $M_n$ ketika $n \to \infty$ dan memberikan hasil yang analog dengan teorema limit sentral untuk maksimum (bukan untuk jumlah), dengan syarat kondisi teknis tertentu pada fungsi distribusi terpenuhi. Jelas bahwa tanpa restriksi lebih lanjut, $M_n$ akan mendekati batas atas dari support distribusi,
\begin{equation}
    \omega = \sup\{x \geq 0 : {}_xq_0 < 1\},
\end{equation}
yang mungkin berhingga atau tak berhingga. Hal ini dapat dilihat dari
\begin{equation}
    P(M_n \leq x) = ({}_xq_0)^n \to \begin{cases}
        0 & \text{jika } x < \omega, \\
        1 & \text{jika } x \geq \omega,
    \end{cases}
\end{equation}
ketika $n \to \infty$.

Namun, setelah $M_n$ dipusatkan dan dinormalisasi secara tepat, distribusinya dapat konvergen ke suatu distribusi limit tertentu. Secara lebih presisi, jika terdapat sekuens bilangan riil $a_n > 0$ dan $b_n \in \mathbb{R}$ sedemikian sehingga sekuens ternormalisasi $(M_n - b_n)/a_n$ konvergen dalam distribusi ke $H$, yaitu
\begin{equation}
    \lim_{n \to \infty} P\left(\frac{M_n - b_n}{a_n} \leq x\right) = \lim_{n \to \infty} \left({}_{{a_n x + b_n}}q_0\right)^n = H(x),
\end{equation}
untuk semua titik kontinuitas dari $H$, maka $H$ adalah distribusi nilai ekstrem tergeneralisasi (\textit{generalized extreme value distribution}, GEV), yaitu $H = H_\xi$ yang diberikan oleh
\begin{equation}
    H_\xi(x) = \begin{cases}
        \exp\left(-(1 + \xi x)_+^{-1/\xi}\right) & \text{jika } \xi \neq 0, \\
        \exp(-\exp(-x)) & \text{jika } \xi = 0,
    \end{cases}
\end{equation}
dengan $y_+ = \max\{y, 0\}$ adalah bagian positif dari $y$ \citep{coles2001introduction, resnick2007heavy}. Domain definisi dari $H_\xi$ adalah $(-1/\xi, +\infty)$ jika $\xi > 0$, $(-\infty, -1/\xi)$ jika $\xi < 0$, dan seluruh garis riil $\mathbb{R}$ jika $\xi = 0$. Parameter $\xi$ yang mengontrol ekor kanan distribusi disebut \textit{tail index} atau \textit{extreme value index}. Tiga distribusi nilai ekstrem klasik adalah kasus khusus dari keluarga GEV: jika $\xi > 0$, diperoleh distribusi Fr\'echet; jika $\xi < 0$, diperoleh distribusi Weibull; dan $\xi = 0$ memberikan distribusi Gumbel.

Dalam aplikasi mortalitas, kasus $\xi > 0$ mengimplikasikan masa hidup dengan ekor berat (\textit{heavy tails}), yang berarti fungsi \textit{hazard} menurun seiring bertambahnya usia. Hal ini bertentangan dengan bukti empiris untuk masa hidup manusia. Dengan demikian, kasus $\xi = 0$ dan $\xi < 0$ yang relevan untuk aplikasi asuransi jiwa \citep{gbari2017extreme}. Perlu dicatat bahwa jika kondisi konvergensi terpenuhi dengan $\xi < 0$, maka $\omega < \infty$, sehingga nilai negatif dari $\xi$ mendukung eksistensi usia maksimum berhingga.

Suatu kondisi cukup untuk konvergensi di atas adalah
\begin{equation}
    \lim_{x \to \omega} \frac{d}{dx}\left(\frac{1}{\mu_x}\right) = \xi,
\end{equation}
dengan $\mu_x$ adalah \textit{force of mortality} pada usia $x$. Secara intuitif, $1/\mu_x$ dapat dipandang sebagai kekuatan resistensi terhadap mortalitas atau kekuatan vitalitas pada usia $x$. Resistensi terhadap mortalitas harus stabil ketika $\xi = 0$ atau menjadi linier secara asimptotik. Nilai $\xi$ negatif mengindikasikan bahwa resistensi pada akhirnya menurun pada usia lanjut. Untuk $\xi < 0$, diperoleh $\omega < \infty$, dan kondisi di atas mengimplikasikan
\begin{equation}
    \lim_{x \to \omega} \left[(\omega - x) \mu_x\right] = -\frac{1}{\xi}.
\end{equation}

\subsection{Pendekatan Peaks Over Threshold (POT)}

Selain pendekatan \textit{block maxima} yang menganalisis nilai maksimum dalam blok-blok data, EVT juga menawarkan pendekatan alternatif yang disebut \textit{Peaks Over Threshold} (POT). Pendekatan ini menganalisis semua nilai observasi yang melebihi suatu ambang batas $u$ yang cukup tinggi, sehingga memanfaatkan informasi dari seluruh data ekstrem yang tersedia, bukan hanya nilai maksimum per blok \citep{coles2001introduction}.

Dalam konteks mortalitas, perhatikan sisa masa hidup $T - x$ pada usia $x$, dengan asumsi $T > x$, yang memiliki fungsi distribusi
\begin{equation}
    {}_sq_x = P(T - x \leq s \mid T > x), \quad s \geq 0.
\end{equation}
Dapat terjadi bahwa untuk usia lanjut $x$ yang besar, distribusi probabilitas kondisional ini menjadi stabil setelah normalisasi, yaitu terdapat fungsi positif $a(\cdot)$ sedemikian sehingga
\begin{equation}
    \lim_{x \to \omega} P\left(\frac{T - x}{a(x)} > s \,\bigg|\, T > x\right) = 1 - G(s), \quad s > 0,
\end{equation}
dengan $G$ adalah fungsi distribusi yang tidak degenerasi. Dapat ditunjukkan bahwa hanya kelas terbatas dari fungsi distribusi yang memenuhi kondisi di atas, yaitu
\begin{equation}
    G(s) = G_\xi(s) = \ln H_\xi(s) = \begin{cases}
        1 - (1 + \xi s)_+^{-1/\xi} & \text{jika } \xi \neq 0, \\
        1 - \exp(-s) & \text{jika } \xi = 0.
    \end{cases}
\end{equation}
Support distribusi ini adalah setengah garis riil positif jika $\xi \geq 0$ dan $[0, -1/\xi]$ jika $\xi < 0$. Keluarga skala terkait yang dikenal sebagai \textit{generalized Pareto distribution} (GPD) didefinisikan sebagai
\begin{equation}
    G_{\xi, \beta}(s) = G_\xi\left(\frac{s}{\beta}\right), \quad \beta > 0,
\end{equation}
dengan $\beta$ adalah parameter skala. Kasus-kasus khusus dari GPD meliputi distribusi Pareto ketika $\xi > 0$, distribusi Pareto tipe II ketika $\xi < 0$, dan distribusi eksponensial negatif ketika $\xi = 0$. Dengan demikian, ketika $\xi = 0$, sisa masa hidup pada usia tinggi menjadi terdistribusi eksponensial negatif secara asimptotik, sehingga \textit{force of mortality} menjadi konstan, sejalan dengan studi empiris yang dilakukan oleh Gampe pada tahun 2010 \citep{gampe2010human}.

Analisis nilai ekstrem untuk maksimum dengan demikian berkaitan erat dengan studi sisa masa hidup. Dapat ditunjukkan bahwa konvergensi ke GEV berlaku jika dan hanya jika konvergensi ke GPD berlaku. Dengan kata lain, $G_\xi$ menggambarkan sisa masa hidup di atas usia yang cukup tua jika dan hanya jika $H_\xi$ mengatur perilaku maksimum sampel, yaitu jika fungsi distribusi $F$ termasuk dalam domain atraksi dari distribusi GEV.

\subsection{Generalized Pareto Distribution dan Teorema Pickands-Balkema-de Haan}

Untuk suatu fungsi $\beta(\cdot)$ yang sesuai, aproksimasi
\begin{equation}
    {}_sq_x \approx G_{\xi; \beta(x)}(s), \quad s \geq 0,
\end{equation}
berlaku untuk $x$ yang cukup besar. Aproksimasi ini dijustifikasi oleh Teorema Pickands-Balkema-de Haan yang menyatakan bahwa
\begin{equation}
    \lim_{x \to \omega} \sup_{s \geq 0} \left|{}_sq_x - G_{\xi, \beta(x)}(s)\right| = 0,
\end{equation}
berlaku dengan syarat $F$ memenuhi kondisi teknis yang cukup umum \citep{balkema1974residual, pickands1975statistical}. Berdasarkan aproksimasi ini, sisa masa hidup pada usia $x$ dapat diperlakukan sebagai sampel acak dari distribusi GPD, dengan syarat $x$ cukup besar.

Jika $\omega < \infty$ (yaitu $\xi < 0$), maka transformasi yang sesuai dari \textit{extreme value index} $\xi$ memiliki interpretasi intuitif. Harapan hidup tersisa pada usia $x$, yang dinotasikan sebagai $e_x$, didefinisikan sebagai
\begin{equation}
    e_x = E[T - x \mid T > x].
\end{equation}
Aarssen dan de Haan pada tahun 1994 menetapkan bahwa untuk $\xi < 0$, sehingga batas atas $\omega$ ada pada rentang masa hidup, konvergensi ke GEV ekuivalen dengan
\begin{equation}
    \lim_{x \to \omega} E\left[\frac{T - x}{\omega - x} \,\bigg|\, T > x\right] = \lim_{x \to \omega} \frac{e_x}{\omega - x} = -\frac{\xi}{1 - \xi} = \alpha.
\end{equation}
Para peneliti ini menyebut $\alpha = \alpha(\xi)$ sebagai \textit{perseverance parameter} dan memberikan penjelasan sebagai berikut \citep{aarssen1994domains}. Perhatikan seorang individu yang masih hidup pada usia lanjut $x$. Rasio $(T - x)/(\omega - x)$ merepresentasikan persentase sisa masa hidup aktual $T - x$ terhadap sisa masa hidup maksimum $\omega - x$. Persentase ini menjadi stabil, secara rata-rata, ketika $x \to \omega$ dan konvergen ke $\alpha$, yang dengan demikian muncul sebagai persentase ekspektasi dari sisa masa hidup maksimum yang mungkin, yang secara efektif digunakan oleh individu tersebut.

Teorema Pickands-Balkema-de Haan memberikan justifikasi teoretis untuk pendekatan POT dalam pemodelan mortalitas usia lanjut. Dalam praktik, untuk menerapkan hasil ini, perlu ditentukan usia ambang $u$ (atau $x^*$) sedemikian sehingga aproksimasi GPD cukup akurat untuk $x \geq u$. Pemilihan usia ambang yang tepat merupakan aspek krusial dalam implementasi model berbasis EVT, dan akan dibahas lebih lanjut dalam konteks konstruksi model \textit{Threshold Life Table} pada sub-bab berikutnya.

\section{Model Mortalitas Dinamis: Model Cairns-Blake-Dowd (CBD)}

Model mortalitas dinamis menangkap perubahan pola mortalitas dari waktu ke waktu, yang penting untuk melakukan proyeksi mortalitas masa depan. Berbeda dengan model statis yang mengasumsikan parameter tetap, model dinamis memodelkan parameter sebagai fungsi waktu atau kohor. Salah satu contoh model mortalitas dinamis adalah Model Cairns-Blake-Dowd (CBD).

Model Cairns-Blake-Dowd (CBD) yang diperkenalkan oleh \citet{cairns2006two} dirancang khusus untuk memodelkan mortalitas pada usia lanjut. Model ini menggunakan struktur yang lebih sederhana dibandingkan Lee-Carter namun tetap efektif untuk aplikasi pada rentang usia terbatas.

Model CBD memodelkan logit dari probabilitas kematian satu tahun $q_{x,t}$ sebagai fungsi linear dari usia:
\begin{equation}
\text{logit}(q_{x,t}) = \ln\left(\frac{q_{x,t}}{1-q_{x,t}}\right) = \kappa_t^{(1)} + \kappa_t^{(2)} (x - \bar{x}),
\label{eq:cbd_basic}
\end{equation}
dimana:
\begin{itemize}
    \item $x$ adalah usia (umumnya untuk rentang usia lanjut, misalnya 60--89 tahun)
    \item $t$ adalah tahun kalender atau indeks waktu
    \item $\bar{x}$ adalah usia rata-rata dalam rentang yang dianalisis (untuk centering)
    \item $\kappa_t^{(1)}$ adalah parameter temporal yang mengatur tingkat mortalitas keseluruhan pada tahun $t$
    \item $\kappa_t^{(2)}$ adalah parameter temporal yang mengatur slope atau gradien mortalitas terhadap usia pada tahun $t$
\end{itemize}

Dari persamaan \eqref{eq:cbd_basic}, probabilitas kematian dapat diperoleh melalui transformasi invers:
\begin{equation}
q_{x,t} = \frac{\exp(\kappa_t^{(1)} + \kappa_t^{(2)} (x - \bar{x}))}{1 + \exp(\kappa_t^{(1)} + \kappa_t^{(2)} (x - \bar{x}))}.
\label{eq:cbd_qx}
\end{equation}

Parameter temporal $\kappa_t^{(1)}$ dan $\kappa_t^{(2)}$ dimodelkan sebagai proses stokastik. Spesifikasi standar menggunakan \textit{random walk with drift}:
\begin{align}
\kappa_t^{(1)} &= \kappa_{t-1}^{(1)} + \mu^{(1)} + \epsilon_t^{(1)}, \label{eq:cbd_rw1} \\
\kappa_t^{(2)} &= \kappa_{t-1}^{(2)} + \mu^{(2)} + \epsilon_t^{(2)}, \label{eq:cbd_rw2}
\end{align}
dimana $\mu^{(1)}$ dan $\mu^{(2)}$ adalah drift parameters, dan $\epsilon_t^{(1)}, \epsilon_t^{(2)}$ adalah error terms yang diasumsikan white noise dengan mean nol.



Dalam penelitian ini, model CBD digunakan sebagai \textit{benchmark} untuk mengevaluasi kemampuan forecasting model Dynamic Smooth Threshold Life Table (DSTLT). Perbandingan dilakukan pada:
\begin{enumerate}
    \item \textbf{In-sample fit}: Kesesuaian model terhadap data historis yang digunakan untuk estimasi
    \item \textbf{Out-of-sample forecast accuracy}: Akurasi proyeksi mortalitas untuk kohor yang tidak termasuk dalam training set
\end{enumerate}

Model CBD dipilih sebagai benchmark karena merupakan model standar dalam literatur aktuaria untuk proyeksi mortalitas usia lanjut dan telah banyak digunakan dalam praktik industri asuransi dan dana pensiun \citep{villegas2018comparative}.

\section{Maximum Likelihood Estimation (MLE)}

Maximum Likelihood Estimation (MLE) merupakan metode estimasi parameter yang paling umum digunakan dalam statistika dan pemodelan mortalitas. Metode ini memilih nilai parameter yang memaksimalkan probabilitas (atau likelihood) dari data yang diamati.

\subsection{Konsep Dasar}

Misalkan terdapat data observasi $\mathbf{x} = (x_1, x_2, \ldots, x_n)$ yang diasumsikan berasal dari distribusi dengan fungsi kepadatan (atau fungsi massa) probabilitas $f(x; \boldsymbol{\theta})$, dimana $\boldsymbol{\theta}$ adalah vektor parameter yang tidak diketahui. Fungsi likelihood didefinisikan sebagai fungsi dari parameter $\boldsymbol{\theta}$ yang menyatakan probabilitas bersama dari data observasi:
\begin{equation}
L(\boldsymbol{\theta}; \mathbf{x}) = \prod_{i=1}^{n} f(x_i; \boldsymbol{\theta}).
\label{eq:likelihood_definition}
\end{equation}

Estimator Maximum Likelihood $\hat{\boldsymbol{\theta}}$ adalah nilai parameter yang memaksimalkan fungsi likelihood:
\begin{equation}
\hat{\boldsymbol{\theta}} = \arg\max_{\boldsymbol{\theta}} L(\boldsymbol{\theta}; \mathbf{x}).
\label{eq:mle_definition}
\end{equation}

Dalam pengerjaannya, lebih mudah untuk mendapatkan solusi dalam bentuk logaritma natural dari likelihood, yang biasa disebut log-likelihood:
\begin{equation}
\ell(\boldsymbol{\theta}; \mathbf{x}) = \ln L(\boldsymbol{\theta}; \mathbf{x}) = \sum_{i=1}^{n} \ln f(x_i; \boldsymbol{\theta}).
\label{eq:loglikelihood_definition}
\end{equation}

Untuk mendapatkan MLE, biasanya dilakukan dengan mencari solusi dari persamaan:
\begin{equation}
\frac{\partial \ell(\boldsymbol{\theta})}{\partial \boldsymbol{\theta}} = \mathbf{0}.
\label{eq:score_equation}
\end{equation}

\subsection{MLE untuk Data Tersensor}

Dalam analisis mortalitas, data sering bersifat tersensor, yang memerlukan penyesuaian dalam konstruksi fungsi likelihood.

Untuk data tersensor interval, dimana usia kematian hanya diketahui berada dalam interval $[x, x+1)$, kontribusi likelihood untuk setiap observasi adalah probabilitas berada dalam interval tersebut:
\begin{equation}
L_i = P(x < X \leq x+1) = F(x+1; \boldsymbol{\theta}) - F(x; \boldsymbol{\theta}).
\label{eq:likelihood_interval}
\end{equation}

Jika terdapat $d_x$ individu yang meninggal dalam interval $[x, x+1)$ dari $l_x$ individu berisiko, dan dengan mengabaikan konstanta kombinatorial, kontribusi log-likelihood adalah:
\begin{equation}
\ell_x = d_x \ln[F(x+1; \boldsymbol{\theta}) - F(x; \boldsymbol{\theta})].
\label{eq:loglik_interval}
\end{equation}

Dalam bentuk fungsi survival:
\begin{equation}
\ell_x = d_x \ln[S(x; \boldsymbol{\theta}) - S(x+1; \boldsymbol{\theta})].
\label{eq:loglik_interval_survival}
\end{equation}

Untuk individu yang tersensor kanan pada usia $\tau$ (masih hidup pada akhir pengamatan), kontribusi likelihood adalah probabilitas bertahan hidup melampaui $\tau$:
\begin{equation}
L_i = P(X > \tau) = S(\tau; \boldsymbol{\theta}).
\label{eq:likelihood_right_censored}
\end{equation}

Jika terdapat $l_\tau$ individu yang tersensor kanan, kontribusi log-likelihood adalah:
\begin{equation}
\ell_{\tau} = l_\tau \ln S(\tau; \boldsymbol{\theta}).
\label{eq:loglik_right_censored}
\end{equation}

Untuk data mortalitas dengan rentang usia dari $x_{\min}$ hingga $\tau$, fungsi log-likelihood total menggabungkan kontribusi dari semua interval usia dan data tersensor kanan:
\begin{equation}
\ell(\boldsymbol{\theta}) = \sum_{x=x_{\min}}^{\tau-1} d_x \ln[S(x; \boldsymbol{\theta}) - S(x+1; \boldsymbol{\theta})] + l_\tau \ln S(\tau; \boldsymbol{\theta}) - l_{x_{\min}} \ln S(x_{\min}; \boldsymbol{\theta}).
\label{eq:loglik_total_censored}
\end{equation}

Normalisasi dengan $S(x_{\min})$ memastikan bahwa likelihood dikondisikan pada bertahan hidup hingga usia awal pengamatan, yang konsisten dengan definisi probabilitas kondisional dalam analisis kohor.

Struktur likelihood ini merupakan dasar untuk estimasi parameter dalam model TLT, STLT, dan DSTLT yang akan dibahas pada Bab 3.



\section{Kriteria Evaluasi Model}

Evaluasi kinerja model mortalitas memerlukan kriteria yang tepat untuk mengukur kesesuaian model terhadap data historis dan akurasi proyeksi untuk data masa depan.

\subsection{Sum of Squared Errors (SSE)}

Sum of Squared Errors mengukur total deviasi kuadrat antara nilai prediksi model dengan nilai observasi:
\begin{equation}
\text{SSE} = \sum_{i=1}^{n} (y_i - \hat{y}_i)^2,
\label{eq:sse}
\end{equation}
dimana $y_i$ adalah nilai observasi (misalnya, $q_x$ empiris atau tingkat mortalitas observasi) dan $\hat{y}_i$ adalah nilai prediksi dari model.

Untuk evaluasi model mortalitas, SSE dapat dihitung berdasarkan:
\begin{itemize}
    \item Probabilitas kematian: $\text{SSE}_q = \sum_x (q_x^{\text{obs}} - q_x^{\text{pred}})^2$
    \item Tingkat mortalitas: $\text{SSE}_\mu = \sum_x (\mu_x^{\text{obs}} - \mu_x^{\text{pred}})^2$
\end{itemize}

Nilai SSE yang lebih kecil mengindikasikan kesesuaian yang lebih baik. Metrik ini terutama berguna untuk model seperti Coale-Kisker yang tidak diestimasi menggunakan MLE, sehingga tidak memiliki nilai log-likelihood yang dapat dibandingkan.

\subsection{In-Sample Fit vs Out-of-Sample Forecast}

Untuk mengevaluasi kemampuan prediksi model, diperlukan ukuran yang membandingkan proyeksi model dengan data aktual yang tidak digunakan dalam estimasi. Evaluasi model dapat dibagi menjadi dua kategori:

\begin{itemize}
    \item \textbf{In-sample fit}: Kesesuaian model terhadap data yang digunakan untuk estimasi parameter (training set). Ukuran ini menilai seberapa baik model menangkap pola dalam data historis.
    
    \item \textbf{Out-of-sample forecast}: Akurasi proyeksi model pada data yang tidak digunakan dalam estimasi (test set). Ukuran ini menilai kemampuan generalisasi dan prediksi model.
\end{itemize}

Untuk model dinamis seperti DSTLT, evaluasi out-of-sample sangat penting karena tujuan utama model adalah untuk melakukan proyeksi mortalitas masa depan.


\chapter{\normalsize MODEL DYNAMIC SMOOTH THRESHOLD LIFE TABLE (DSTLT)}

Pemodelan mortalitas pada usia lanjut ekstrem menghadapi tantangan yang tidak dapat diatasi sepenuhnya oleh model-model mortalitas tradisional. Keterbatasan data pada usia sangat lanjut, volatilitas tinggi dari estimasi empiris, dan kebutuhan untuk menentukan batas akhir tabel mortalitas secara objektif mendorong pengembangan metode yang lebih canggih. Bab ini membahas mengenai perkembangan bertahap dari pemodelan mortalitas usia lanjut, dimulai dari model \textit{Threshold Life Table} (TLT) sebagai model dasar, kemudian pengembangan \textit{Smoothed Threshold Life Table} (STLT) sebagai perbaikan, hingga \textit{Dynamic Smooth Threshold Life Table} (DSTLT) sebagai pengembangan untuk keperluan peramalan.

Model TLT yang diperkenalkan oleh Li et al. (2008) mengatasi masalah terkait penentuan usia penutupan tabel mortalitas yang seringkali bersifat sembarang. Model TLT memiliki struktur \textit{piecewise} yang menggabungkan distribusi Gompertz untuk usia non-ekstrem dengan \textit{Generalized Pareto Distribution} (GPD) untuk usia lanjut ekstrem. Namun, model TLT memiliki keterbatasan terkait kontinuitas fungsi \textit{hazard} pada titik transisi antara kedua bagian model. Ketidakmulusan fungsi \textit{hazard} ini menimbulkan implikasi yang tidak realistis secara biologis (nanti dimasukin sumbernya), dimana laju mortalitas sesaat dapat mengalami loncatan mendadak pada usia ambang tertentu. Keterbatasan inilah yang memotivasi pengembangan model STLT oleh Huang et al. (2020), yang memperkenalkan \textit{smoothing constraint} untuk memastikan kontinuitas fungsi \textit{hazard}.

Walaupun model STLT berhasil mengatasi masalah kontinuitas, terdapat keterbatasan dalam peramalan mortalitas. Pengamatan empiris terhadap parameter model STLT antar kohor menunjukkan adanya pola. Hal ini mendorong pengembangan model DSTLT, yang mengintegrasikan komponen dinamis untuk menangkap perubahan mortalitas antar cohort.

Bab ini akan menganalisis secara mendalam setiap tahap pengembangan model, dimulai dari formula dasar TLT, penurunan rumus \textit{smoothing constraint} pada STLT, hingga konstruksi komponen dinamis pada DSTLT.

\section{Model Threshold Life Table (TLT)}

Model \textit{Threshold Life Table} (TLT) yang diperkenalkan oleh \citet{li2008threshold} menggunakan pendekatan \textit{piecewise} yang membagi distribusi usia kematian menjadi dua bagian berdasarkan usia ambang (\textit{threshold age}) $u$. Model ini dikembangkan untuk mengatasi keterbatasan model mortalitas tradisional dalam menangani karakteristik yang berbeda antara usia lanjut non-ekstrem dan usia lanjut ekstrem.

\subsection{Model Threshold Life Table (TLT)}

Misalkan $X$ adalah variabel acak yang menyatakan usia kematian seseorang. Model TLT mendefinisikan fungsi distribusi kumulatif (CDF) $F(x) = P(X \leq x)$ secara \textit{piecewise} sebagai berikut:

\paragraph{Komponen Gompertz ($x \leq N$):}
Untuk usia di bawah atau sama dengan usia ambang $u$, digunakan distribusi Gompertz:
\begin{equation}
F(x) = 1 - \exp\left(-\frac{B}{\ln C}(C^x - 1)\right), \quad x \leq N,
\label{eq:tlt_gompertz_cdf}
\end{equation}
dengan parameter $B > 0$ sebagai parameter skala dan $C > 1$ sebagai parameter bentuk yang menentukan laju peningkatan mortalitas seiring bertambahnya usia. Fungsi survival untuk bagian ini adalah:
\begin{equation}
S(x) = \exp\left(-\frac{B}{\ln C}(C^x - 1)\right), \quad x \leq N.
\label{eq:tlt_gompertz_survival}
\end{equation}

Fungsi kepadatan probabilitas (PDF) diperoleh melalui diferensiasi $f(x) = -S'(x)$:
\begin{equation}
f(x) = BC^x \exp\left(-\frac{B}{\ln C}(C^x - 1)\right), \quad x \leq N.
\label{eq:tlt_gompertz_pdf}
\end{equation}

Fungsi hazard untuk bagian Gompertz adalah:
\begin{equation}
h(x) = \frac{f(x)}{S(x)} = BC^x, \quad x \leq N,
\label{eq:tlt_gompertz_hazard}
\end{equation}
yang menunjukkan peningkatan eksponensial seiring bertambahnya usia.

\paragraph{Komponen Generalized Pareto Distribution ($x > u$):}
Untuk usia di atas usia ambang $u$, digunakan Generalized Pareto Distribution (GPD) yang merupakan bagian dari Extreme Value Theory. GPD didefinisikan secara kondisional untuk \textit{exceedances} di atas ambang $u$. Fungsi survival kondisional untuk bagian ini adalah:
\begin{equation}
G(x-u; \theta, \gamma) = \left(1 + \frac{\gamma(x-u)}{\theta}\right)^{-1/\gamma}, \quad x > u,
\label{eq:gpd_survival_conditional}
\end{equation}
dengan $\theta > 0$ adalah parameter skala dan $\gamma$ adalah parameter bentuk yang menentukan perilaku ekor distribusi. Fungsi survival keseluruhan untuk $x > u$ adalah:
\begin{equation}
S(x) = S(u) \cdot G(x-u; \theta, \gamma), \quad x > u,
\label{eq:tlt_gpd_survival}
\end{equation}
yang memastikan kontinuitas fungsi survival pada $x = u$.

Fungsi kepadatan probabilitas kondisional untuk bagian GPD adalah:
\begin{equation}
g(x-u; \theta, \gamma) = \frac{1}{\theta}\left(1 + \frac{\gamma(x-u)}{\theta}\right)^{-(1+1/\gamma)}, \quad x > u.
\label{eq:gpd_pdf_conditional}
\end{equation}

\subsubsection{Parameter Bentuk GPD dan Perilaku Ekor}

Parameter bentuk $\gamma$ dalam distribusi GPD menentukan karakteristik ekor distribusi:

\begin{itemize}
    \item \textbf{$\gamma < 0$}: Distribusi memiliki batas atas (\textit{highest attained age}) $\omega = u + \theta/|\gamma|$. Kondisi ini mengimplikasikan bahwa terdapat usia maksimum yang dapat dicapai oleh populasi.
    
    \item \textbf{$\gamma = 0$}: Distribusi ekivalen dengan distribusi eksponensial, yang memiliki ekor yang menurun secara eksponensial.
    
    \item \textbf{$\gamma > 0$}: Distribusi memiliki ekor yang sangat berat dan tidak memiliki batas atas. Kondisi ini mengimplikasikan tidak ada batas teoretis untuk usia maksimum.
\end{itemize}

Dalam konteks mortalitas manusia, nilai $\gamma < 0$ umumnya lebih masuk akal secara biologis karena terdapat batas praktis untuk umur manusia \citep{dong2016evidence}.

\subsection{Konstruksi Fungsi Likelihood}

Salah satu karakteristik penting dari model TLT adalah kemampuan untuk melakukan estimasi parameter secara terpisah antara bagian Gompertz dan bagian GPD. Untuk memahami mengapa hal ini dimungkinkan, perlu dilakukan analisis terhadap struktur fungsi likelihood model TLT.

\subsubsection{Struktur Data Mortalitas}

Data mortalitas yang digunakan dalam estimasi model TLT umumnya berbentuk data agregat kohor dengan struktur sebagai berikut:

\begin{itemize}
    \item $d_x$: jumlah kematian yang terjadi antara usia $x$ dan $x+1$
    \item $l_x$: jumlah individu yang masih hidup pada awal usia $x$ (\textit{exposure})
    \item $x_{\text{mulai}}$: usia awal pengamatan (biasanya 65 tahun)
    \item $\tau$: usia maksimum pengamatan dalam data
\end{itemize}

Data seperti ini bersifat tersensor interval (\textit{interval-censored}) karena usia kematian dicatat dalam satuan tahun, dan tersensor kanan (\textit{right-censored}) karena beberapa individu mungkin masih hidup pada akhir periode pengamatan.

\subsection{Konstruksi dan Dekomposisi Fungsi Likelihood}

Salah satu karakteristik penting dari model TLT adalah kemampuan untuk melakukan estimasi parameter secara terpisah antara bagian Gompertz dan bagian GPD. Untuk memahami properti ini, perlu dilakukan analisis terhadap struktur fungsi likelihood model TLT.

\subsubsection{Struktur Data Mortalitas}

Data mortalitas yang digunakan dalam estimasi model TLT umumnya berbentuk data agregat kohor dengan struktur sebagai berikut:

\begin{itemize}
    \item $d_x$: jumlah kematian yang terjadi antara usia $x$ dan $x+1$
    \item $l_x$: jumlah individu yang masih hidup pada awal usia $x$ (\textit{exposure})
    \item $x_{\text{mulai}}$: usia awal pengamatan (biasanya 65 tahun)
    \item $\tau$: usia maksimum pengamatan dalam data
\end{itemize}

Data seperti ini bersifat tersensor interval (\textit{interval-censored}) karena usia kematian dicatat dalam satuan tahun, dan tersensor kanan (\textit{right-censored}) karena beberapa individu mungkin masih hidup pada akhir periode pengamatan.

\subsubsection{Fungsi Likelihood Model TLT}

Fungsi likelihood untuk model TLT dengan parameter $\boldsymbol{\theta} = (B, C, \theta, \gamma, N)$ dapat langsung dikonstruksi berdasarkan probabilitas kematian pada setiap interval usia dan probabilitas bertahan hidup melampaui usia maksimum pengamatan. Untuk setiap usia $x$, probabilitas kematian antara $x$ dan $x+1$ relatif terhadap probabilitas bertahan hidup hingga usia awal pengamatan adalah:
\begin{equation}
\frac{S(x) - S(x+1)}{S(x_{\text{mulai}})},
\label{eq:prob_death_relative}
\end{equation}
yang menyatakan proporsi populasi awal yang meninggal pada interval usia $[x, x+1)$.

Dengan mengabaikan konstanta kombinatorial yang tidak mempengaruhi estimator Maximum Likelihood, fungsi likelihood untuk model TLT dapat ditulis sebagai:
\begin{equation}
\begin{split}
L(B, C, \theta, \gamma; u) = &\prod_{x=x_{\text{mulai}}}^{N-1} \left(\frac{S(x) - S(x+1)}{S(x_{\text{mulai}})}\right)^{d_x} \\
&\times \prod_{x=u}^{\tau-1} \left(\frac{S(x) - S(x+1)}{S(x_{\text{mulai}})}\right)^{d_x} \\
&\times \left(\frac{S(\tau)}{S(x_{\text{mulai}})}\right)^{l_\tau},
\end{split}
\label{eq:tlt_likelihood_full}
\end{equation}
dimana fungsi survival $S(x)$ didefinisikan secara \textit{piecewise}:
\begin{equation}
S(x) = 
\begin{cases}
\exp\left(-\frac{B}{\ln C}(C^x - 1)\right), & x \leq N \\[10pt]
\exp\left(-\frac{B}{\ln C}(C^u - 1)\right) \cdot \left(1 + \frac{\gamma(x-u)}{\theta}\right)^{-1/\gamma}, & x > u.
\end{cases}
\label{eq:tlt_survival_piecewise}
\end{equation}

Suku pertama dalam persamaan \eqref{eq:tlt_likelihood_full} merepresentasikan kontribusi kematian pada usia non-ekstrem ($x < u$), suku kedua merepresentasikan kontribusi kematian pada usia ekstrem ($x \geq N$), dan suku ketiga merepresentasikan kontribusi dari individu yang bertahan hidup melampaui usia maksimum pengamatan (tersensor kanan).

\subsubsection{Dekomposisi Log-Likelihood}

Dengan mengambil logaritma dari persamaan \eqref{eq:tlt_likelihood_full}, diperoleh fungsi log-likelihood:
\begin{equation}
\begin{split}
\ell(B, C, \theta, \gamma; u) = &\sum_{x=x_{\text{mulai}}}^{N-1} d_x \ln(S(x) - S(x+1)) + \sum_{x=u}^{\tau-1} d_x \ln(S(x) - S(x+1)) \\
&+ l_\tau \ln(S(\tau)) - l_{x_{\text{mulai}}} \ln(S(x_{\text{mulai}})).
\end{split}
\label{eq:tlt_loglik_expanded}
\end{equation}

Kunci dari dekomposisi terletak pada fakta bahwa tidak ada parameter yang dibagi antara bagian Gompertz dan bagian GPD. Bagian Gompertz hanya melibatkan parameter $\{B, C\}$, sedangkan bagian GPD hanya melibatkan parameter $\{\theta, \gamma\}$. Dengan memisahkan suku-suku yang hanya bergantung pada masing-masing komponen, fungsi log-likelihood dapat ditulis sebagai jumlah dua fungsi independen:

\begin{equation}
\ell(B, C, \theta, \gamma; u) = \ell_1(B, C; u) + \ell_2(\theta, \gamma; u),
\label{eq:tlt_decomposition}
\end{equation}

dimana komponen pertama (bagian Gompertz):
\begin{equation}
\ell_1(B, C; u) = \sum_{x=x_{\text{mulai}}}^{N-1} d_x \ln(S(x) - S(x+1)) + l_u \ln(S(u)) - l_{x_{\text{mulai}}} \ln(S(x_{\text{mulai}})),
\label{eq:tlt_gompertz_loglik}
\end{equation}
dengan $S(x) = \exp\left(-\frac{B}{\ln C}(C^x - 1)\right)$ untuk $x \leq N$, dan komponen kedua (bagian GPD):
\begin{equation}
\ell_2(\theta, \gamma; u) = \sum_{x=u}^{\tau-1} d_x \ln(S(x) - S(x+1)) + l_\tau \ln(S(\tau)) - l_u \ln(S(u)),
\label{eq:tlt_gpd_loglik}
\end{equation}
dengan $S(x) = S(u) \cdot \left(1 + \frac{\gamma(x-u)}{\theta}\right)^{-1/\gamma}$ untuk $x > u$.

Perhatikan bahwa $\ell_1$ hanya bergantung pada parameter $(B, C)$ dan $\ell_2$ hanya bergantung pada parameter $(\theta, \gamma)$, sementara suku $l_u \ln(S(u))$ muncul dengan tanda berlawanan pada kedua komponen sehingga saling menghilangkan dalam penjumlahan, menghasilkan bentuk \eqref{eq:tlt_loglik_expanded}.

\subsubsection{Implikasi untuk Estimasi Parameter}

Struktur aditif pada persamaan \eqref{eq:tlt_decomposition} memiliki implikasi penting untuk prosedur estimasi. Untuk nilai $u$ yang tetap, estimasi parameter dapat dilakukan secara terpisah:

\begin{enumerate}
    \item \textbf{Estimasi parameter Gompertz}: 
    \begin{equation}
    \{\hat{B}(u), \hat{C}(u)\} = \arg\max_{B,C} \ell_1(B, C; u),
    \label{eq:mle_gompertz}
    \end{equation}
    yang hanya melibatkan optimisasi dua parameter pada data usia $x_{\text{mulai}}$ hingga $u$.
    
    \item \textbf{Estimasi parameter GPD}: 
    \begin{equation}
    \{\hat{\theta}(u), \hat{\gamma}(u)\} = \arg\max_{\theta,\gamma} \ell_2(\theta, \gamma; u),
    \label{eq:mle_gpd}
    \end{equation}
    yang hanya melibatkan optimisasi dua parameter pada data usia $u$ hingga $\tau$.
\end{enumerate}

\subsection{Pemilihan Usia Ambang Optimal}

Meskipun parameter untuk bagian Gompertz dan GPD dapat diestimasi secara terpisah, pemilihan usia ambang $u$ tetap memerlukan pendekatan yang mempertimbangkan kesesuaian keseluruhan model. \citet{li2008threshold} mengembangkan metode berbasis \textit{profile likelihood} untuk pemilihan usia ambang optimal.

\subsubsection{Konsep Profile Likelihood}

\textit{Profile likelihood} adalah fungsi likelihood yang dievaluasi pada estimator Maximum Likelihood (MLE) dari parameter lain, dipandang sebagai fungsi dari parameter yang tersisa. Dalam konteks model TLT, untuk setiap kandidat nilai $u$, dihitung:

\begin{equation}
\ell_p(u) = \ell_1(\hat{B}(u), \hat{C}(u); u) + \ell_2(\hat{\theta}(u), \hat{\gamma}(u); u),
\label{eq:profile_likelihood}
\end{equation}

dimana $\hat{B}(u), \hat{C}(u), \hat{\theta}(u), \hat{\gamma}(u)$ adalah estimator MLE yang bergantung pada nilai $u$.

\subsubsection{Algoritma Pemilihan Usia Ambang}

Berdasarkan implementasi \citet{li2008threshold}, algoritma pemilihan usia ambang dilakukan secara sistematik dengan langkah-langkah berikut:

\begin{enumerate}
    \item \textbf{Inisialisasi}:
    \begin{itemize}
        \item Tentukan rentang kandidat usia ambang: $u \in \{u_{\min}, u_{\min} + 1, \ldots, u_{\max}\}$
        \item Umumnya, $u_{\min} = 85$ dan $u_{\max} = 98$ atau 100
        \item Set $\ell_p^* = -\infty$ dan $\hat{u} = \emptyset$
    \end{itemize}
    
    \item \textbf{Iterasi untuk setiap kandidat $u$}:
    \begin{enumerate}
        \item Estimasi parameter Gompertz $\{\hat{B}(u), \hat{C}(u)\}$ dengan memaksimalkan $\ell_1(B, C; u)$ menggunakan algoritma optimisasi numerik (misalnya, \texttt{optim()} dalam R dengan metode L-BFGS-B)
        
        \item Estimasi parameter GPD $\{\hat{\theta}(u), \hat{\gamma}(u)\}$ dengan memaksimalkan $\ell_2(\theta, \gamma; u)$
        
        \item Hitung profile log-likelihood: $\ell_p(u)$ menggunakan persamaan \eqref{eq:profile_likelihood}
        
        \item Jika $\ell_p(u) > \ell_p^*$, update $\ell_p^* = \ell_p(u)$ dan $\hat{u} = u$
    \end{enumerate}
    
    \item \textbf{Output}: Usia ambang optimal $\hat{u}$ dan parameter model terkait $\{\hat{B}(\hat{u}), \hat{C}(\hat{u}), \hat{\theta}(\hat{u}), \hat{\gamma}(\hat{u})\}$
\end{enumerate}

Pendekatan \textit{profile likelihood} ini memberikan metode objektif untuk menentukan titik transisi yang optimal antara model Gompertz dan GPD berdasarkan kesesuaian data secara keseluruhan.

\subsection{Keterbatasan Model TLT}

Meskipun model TLT memberikan pendekatan yang lebih objektif dan berbasis data untuk pemodelan mortalitas usia lanjut, Model TLT tidak menjamin kemulusan transisi fungsi hazard pada usia ambang $u$. Secara umum, $\lim_{x \to u^-} h(x) \neq \lim_{x \to u^+} h(x)$, yang dapat menimbulkan loncatan mendadak dalam laju mortalitas sesaat. Kondisi ini kurang realistis secara biologis karena mengimplikasikan perubahan mortalitas yang tiba-tiba. Hal ini menjadi motivasi utama pengembangan model STLT yang akan dibahas pada bagian selanjutnya, dimana \textit{smoothing constraint} diperkenalkan untuk memastikan kemulusan fungsi hazard pada usia ambang.

\section{Model Smoothed Threshold Life Table (STLT)}

Model Threshold Life Table (TLT) yang telah dibahas pada bagian sebelumnya memberikan kerangka kerja yang solid untuk pemodelan mortalitas usia lanjut dengan mengintegrasikan Extreme Value Theory. Namun, model ini memiliki keterbatasan mendasar: fungsi hazard tidak dijamin kontinu pada usia ambang $u$. Bagian ini membahas pengembangan model Smoothed Threshold Life Table (STLT) yang mengatasi keterbatasan tersebut melalui penambahan \textit{smoothing constraint}.

\subsection{Masalah Diskontinuitas Fungsi Hazard pada Model TLT}

\subsubsection{Fungsi Hazard Model TLT}

Fungsi hazard merepresentasikan laju mortalitas sesaat pada usia tertentu, kondisional pada bertahan hidup hingga usia tersebut. Untuk model TLT, fungsi hazard didefinisikan secara \textit{piecewise}:

\begin{equation}
h(x) = 
\begin{cases}
BC^x, & x \leq N \quad \text{(Gompertz)} \\[8pt]
\displaystyle\frac{1}{\theta + \gamma(x-u)}, & x > u \quad \text{(GPD)}
\end{cases}
\label{eq:tlt_hazard_piecewise}
\end{equation}

Untuk bagian Gompertz ($x \leq N$), fungsi hazard diperoleh dari:
\begin{equation}
h_1(x) = \frac{f(x)}{S(x)} = \frac{BC^x \exp\left(-\frac{B}{\ln C}(C^x - 1)\right)}{\exp\left(-\frac{B}{\ln C}(C^x - 1)\right)} = BC^x,
\label{eq:hazard_gompertz_derivation}
\end{equation}
yang menunjukkan pertumbuhan eksponensial seiring bertambahnya usia.

Untuk bagian GPD ($x > u$), fungsi hazard diperoleh dari fungsi survival kondisional $G(x-u; \theta, \gamma) = \left(1 + \frac{\gamma(x-u)}{\theta}\right)^{-1/\gamma}$. Dengan diferensiasi dan manipulasi aljabar:
\begin{align}
h_2(x) &= \frac{g(x-u; \theta, \gamma)}{G(x-u; \theta, \gamma)} \nonumber \\
&= \frac{\frac{1}{\theta}\left(1 + \frac{\gamma(x-u)}{\theta}\right)^{-(1+1/\gamma)}}{\left(1 + \frac{\gamma(x-u)}{\theta}\right)^{-1/\gamma}} \nonumber \\
&= \frac{1}{\theta + \gamma(x-u)}.
\label{eq:hazard_gpd_derivation}
\end{align}

\subsubsection{Evaluasi Kontinuitas pada Usia Ambang}

Pada titik transisi $x = u$, nilai fungsi hazard dari kedua komponen adalah:
\begin{itemize}
    \item Limit dari kiri (Gompertz): $\lim_{x \to u^-} h(x) = BC^u$
    \item Limit dari kanan (GPD): $\lim_{x \to u^+} h(x) = \frac{1}{\theta + \gamma(N-N)} = \frac{1}{\theta}$
\end{itemize}

Secara umum, tidak terdapat jaminan bahwa $BC^u = \frac{1}{\theta}$. Dalam estimasi model TLT, parameter $(B, C)$ dan $(\theta, \gamma)$ diestimasi secara independen untuk masing-masing komponen, sehingga kondisi kontinuitas tidak secara otomatis terpenuhi. Ketidaksamaan ini mengakibatkan diskontinuitas (loncatan) pada fungsi hazard:
\begin{equation}
\Delta h(u) = h_1(u) - h_2(u) = BC^u - \frac{1}{\theta} \neq 0.
\label{eq:hazard_jump}
\end{equation}



Loncatan mendadak dalam laju mortalitas sesaat tidak memiliki justifikasi biologis. Proses penuaan dan deteriorasi fisiologis bersifat gradual dan kontinu, sehingga laju mortalitas sesaat seharusnya juga berubah secara mulus. Keterbatasan ini memotivasi pengembangan model STLT yang memastikan kemulusan fungsi hazard melalui penambahan \textit{constraint} matematis.

\subsection{Derivasi Smoothing Constraint}

Model STLT dikembangkan dengan menambahkan kondisi kontinuitas eksplisit pada fungsi hazard di usia ambang $u$. Bagian ini menyajikan derivasi matematis lengkap untuk \textit{smoothing constraint}.

Untuk memastikan fungsi hazard kontinu pada $x = u$, dikenakan kondisi:
\begin{equation}
h_1(u) = h_2(u).
\label{eq:smoothing_condition}
\end{equation}

Substitusi definisi fungsi hazard dari persamaan \eqref{eq:tlt_hazard_piecewise}:
\begin{equation}
BC^u = \frac{1}{\theta + \gamma(N-N)}.
\label{eq:smoothing_substitution}
\end{equation}

Perhatikan bahwa pada $x = u$, suku $\gamma(x-u)$ bernilai nol, sehingga persamaan menjadi:
\begin{equation}
BC^u = \frac{1}{\theta}.
\label{eq:smoothing_simplified}
\end{equation}

Dari persamaan \eqref{eq:smoothing_simplified}, diperoleh hubungan eksplisit antara parameter skala GPD dengan parameter Gompertz:
\begin{equation}
\boxed{\theta = \frac{1}{BC^u}}
\label{eq:theta_constraint}
\end{equation}

Persamaan \eqref{eq:theta_constraint} merupakan hasil kunci dari \textit{smoothing constraint}. Hubungan ini menunjukkan bahwa parameter skala GPD ($\theta$) bukan lagi parameter bebas yang dapat diestimasi secara independen, melainkan merupakan fungsi deterministik dari parameter Gompertz $(B, C)$ dan usia ambang $u$.

\subsection{Formula Model Smooth Threshold Life Table (STLT)}

Dengan menerapkan \textit{smoothing constraint}, model STLT dapat diformulasikan secara lengkap dengan parameter efektif $(B, C, \gamma, N)$.

\\subsection{Formulasi Model STLT}

Dengan menerapkan \textit{smoothing constraint} $\theta = \frac{1}{BC^u}$, model STLT dapat diformulasikan secara lengkap. Berbeda dengan model TLT, parameter skala GPD ($\theta$) bukan lagi parameter bebas, melainkan ditentukan oleh parameter Gompertz dan usia ambang.

\subsubsection{Komponen Gompertz ($x \leq N$)}

Untuk usia yang tidak melebihi usia ambang, model STLT menggunakan distribusi Gompertz yang identik dengan model TLT:

\paragraph{Fungsi Distribusi Kumulatif:}
\begin{equation}
F(x) = 1 - \exp\left(-\frac{B}{\ln C}(C^x - 1)\right), \quad x \leq N.
\label{eq:stlt_gompertz_cdf}
\end{equation}

\paragraph{Fungsi Survival:}
\begin{equation}
S(x) = \exp\left(-\frac{B}{\ln C}(C^x - 1)\right), \quad x \leq N.
\label{eq:stlt_gompertz_survival}
\end{equation}

\paragraph{Fungsi Hazard:}
\begin{equation}
h(x) = BC^x, \quad x \leq N.
\label{eq:stlt_gompertz_hazard}
\end{equation}

Parameter model untuk bagian ini adalah $(B, C)$ dengan interpretasi yang sama seperti pada model TLT: $B > 0$ adalah parameter skala yang merepresentasikan tingkat mortalitas dasar, dan $C > 1$ adalah parameter bentuk yang mengatur laju peningkatan mortalitas.

\subsubsection{Komponen GPD ($x > u$)}

Untuk usia yang melebihi usia ambang, model STLT menggunakan Generalized Pareto Distribution dengan parameter skala yang ditentukan oleh \textit{smoothing constraint}.

\paragraph{Fungsi Distribusi Kumulatif:}

Dengan substitusi $\theta = \frac{1}{BC^u}$ ke dalam formulasi GPD, diperoleh tiga kasus bergantung pada nilai parameter bentuk $\gamma$:

\textbf{Kasus 1: $\gamma > 0$ (ekor berat, tidak ada batas atas usia)}
\begin{equation}
F(x) = 1 - S(u) \left(1 + \gamma \cdot BC^u(x-u)\right)^{-1/\gamma}, \quad x > u.
\label{eq:stlt_gpd_cdf_positive}
\end{equation}

\textbf{Kasus 2: $\gamma = 0$ (ekor eksponensial, tidak ada batas atas usia)}

Untuk $\gamma \to 0$, GPD konvergen ke distribusi eksponensial:
\begin{equation}
F(x) = 1 - S(u) \exp\left(-BC^u(x-u)\right), \quad x > u.
\label{eq:stlt_gpd_cdf_zero}
\end{equation}

\textbf{Kasus 3: $\gamma < 0$ (ekor terbatas, terdapat batas atas usia)}
\begin{equation}
F(x) = 1 - S(u) \left(1 - |\gamma| \cdot BC^u(x-u)\right)^{1/|\gamma|}, \quad N < x < u + \frac{1}{BC^u|\gamma|}.
\label{eq:stlt_gpd_cdf_negative}
\end{equation}

Untuk kasus $\gamma < 0$, terdapat batas atas usia (\textit{highest attained age}):
\begin{equation}
\omega = u + \frac{\theta}{|\gamma|} = u + \frac{1}{BC^u|\gamma|},
\label{eq:stlt_omega}
\end{equation}
dimana $F(\omega) = 1$ dan $S(\omega) = 0$.

\paragraph{Fungsi Survival:}

Fungsi survival untuk bagian GPD dapat ditulis sebagai:

\textbf{Untuk $\gamma \neq 0$:}
\begin{equation}
S(x) = S(u) \left(1 + \gamma \cdot BC^u(x-u)\right)^{-1/\gamma}, \quad x > u,
\label{eq:stlt_gpd_survival_general}
\end{equation}
dengan batasan $x < \omega$ jika $\gamma < 0$.

\textbf{Untuk $\gamma = 0$:}
\begin{equation}
S(x) = S(u) \exp\left(-BC^u(x-u)\right), \quad x > u.
\label{eq:stlt_gpd_survival_zero}
\end{equation}

Perhatikan bahwa dalam semua formulasi di atas, nilai $S(u)$ dihitung menggunakan fungsi survival Gompertz pada persamaan \eqref{eq:stlt_gompertz_survival}, sehingga memastikan kontinuitas fungsi survival pada $x = u$.

\paragraph{Fungsi Hazard:}

Fungsi hazard untuk bagian GPD adalah:
\begin{equation}
h(x) = \frac{1}{\frac{1}{BC^u} + \gamma(x-u)} = \frac{BC^u}{1 + \gamma \cdot BC^u(x-u)}, \quad x > u.
\label{eq:stlt_gpd_hazard}
\end{equation}

\subsubsection{Verifikasi Kontinuitas Fungsi Hazard}

Keunggulan utama model STLT adalah kontinuitas fungsi hazard pada usia ambang $u$. Hal ini dapat diverifikasi dengan mengevaluasi limit dari kiri dan kanan:

\textbf{Limit dari kiri (Gompertz):}
\begin{equation}
\lim_{x \to u^-} h(x) = BC^u.
\label{eq:hazard_left_limit}
\end{equation}

\textbf{Limit dari kanan (GPD):}
\begin{equation}
\lim_{x \to u^+} h(x) = \frac{BC^u}{1 + \gamma \cdot BC^u(N-N)} = \frac{BC^u}{1 + 0} = BC^u.
\label{eq:hazard_right_limit}
\end{equation}

Karena kedua limit bernilai sama, maka:
\begin{equation}
h(u^-) = h(u^+) = BC^u,
\label{eq:hazard_continuity}
\end{equation}
yang membuktikan bahwa fungsi hazard kontinu pada $x = u$. Kontinuitas ini memastikan tidak ada loncatan mendadak dalam laju mortalitas sesaat, sehingga menghasilkan transisi yang lebih realistis secara biologis antara usia lanjut non-ekstrem dan usia lanjut ekstrem.

\subsubsection{Interpretasi Parameter}

Model STLT memiliki tiga parameter bebas untuk estimasi: $(B, C, \gamma)$, dengan parameter keempat $\theta$ ditentukan oleh constraint. Interpretasi masing-masing parameter adalah:

\begin{itemize}
    \item \textbf{$B$}: Parameter skala Gompertz yang mengatur tingkat mortalitas dasar. Nilai $B$ yang lebih besar mengindikasikan tingkat mortalitas yang lebih tinggi pada usia awal periode pengamatan.
    
    \item \textbf{$C$}: Parameter bentuk Gompertz yang mengatur laju akselerasi mortalitas. Nilai $C$ yang lebih besar mengindikasikan percepatan mortalitas yang lebih cepat seiring bertambahnya usia.
    
    \item \textbf{$\gamma$}: Parameter bentuk GPD yang menentukan karakteristik ekor distribusi dan keberadaan batas atas usia. Nilai $\gamma < 0$ (yang umum dalam aplikasi mortalitas manusia) mengimplikasikan adanya batas maksimum usia yang dapat dicapai.
    
    \item \textbf{$\theta = \frac{1}{BC^u}$}: Parameter skala GPD yang secara otomatis ditentukan oleh constraint untuk memastikan kontinuitas hazard. Parameter ini mengatur dispersi distribusi exceedance di atas usia ambang.
\end{itemize}

Hubungan fungsional antara $\theta$ dan parameter lainnya menciptakan struktur model yang terintegrasi, dimana kedua komponen (Gompertz dan GPD) tidak lagi independen tetapi saling terhubung melalui kondisi kemulusan.

\subsection{Struktur Likelihood Model STLT}

Fungsi likelihood untuk model STLT dengan parameter $\boldsymbol{\theta} = (B, C, \gamma, N)$ adalah:
\begin{equation}
\begin{split}
L_{\text{STLT}}(B, C, \gamma; u) = &\prod_{x=x_{\text{mulai}}}^{N-1} \left(\frac{S(x) - S(x+1)}{S(x_{\text{mulai}})}\right)^{d_x} \\
&\times \prod_{x=u}^{\tau-1} \left(\frac{S(x) - S(x+1)}{S(x_{\text{mulai}})}\right)^{d_x} \\
&\times \left(\frac{S(\tau)}{S(x_{\text{mulai}})}\right)^{l_\tau},
\end{split}
\label{eq:stlt_likelihood}
\end{equation}
dengan fungsi survival $S(x)$ yang didefinisikan pada persamaan \eqref{eq:stlt_survival}.

Fungsi log-likelihood adalah:
\begin{equation}
\ell(B, C, \gamma; u) = \sum_{x=x_{\text{mulai}}}^{\tau-1} d_x \ln(S(x) - S(x+1)) + l_\tau \ln(S(\tau)) - l_{x_{\text{mulai}}} \ln(S(x_{\text{mulai}})).
\label{eq:stlt_loglikelihood}
\end{equation}

Struktur likelihood ini menyebabkan tidakmungkin dekomposisi likelihood seperti model TLT.

\subsection{Estimasi Parameter Model STLT}

Estimasi parameter model STLT menggunakan metode Maximum Likelihood Estimation (MLE) dengan prosedur \textit{profile likelihood} untuk pemilihan usia ambang optimal. Algoritma estimasi dapat dirangkum sebagai berikut:

\begin{enumerate}
    \item \textbf{Inisialisasi}:
    \begin{itemize}
        \item Tentukan rentang kandidat usia ambang: $u \in \{u_{\min}, u_{\min}+1, \ldots, u_{\max}\}$
        \item Umumnya $u_{\min} = 85$ dan $u_{\max} = 100$
        \item Set $\ell_p^* = -\infty$ dan $\hat{u} = \emptyset$
    \end{itemize}
    
    \item \textbf{Iterasi untuk setiap kandidat $u$}:
    \begin{enumerate}
        \item Definisikan fungsi objektif:
        \begin{equation}
        \ell(B, C, \gamma; u) = \sum_{x=x_{\text{mulai}}}^{\tau-1} d_x \ln(S(x) - S(x+1)) + l_\tau \ln(S(\tau)) - l_{x_{\text{mulai}}} \ln(S(x_{\text{mulai}}))
        \end{equation}
        dengan $S(x)$ menggunakan constraint $\theta = \frac{1}{BC^u}$.
        
        \item Estimasi parameter $(B, C, \gamma)$ secara simultan menggunakan algoritma optimisasi numerik (misalnya, L-BFGS-B dalam R) dengan batasan $B > 0$, $C > 1$:
        \begin{equation}
        \{\hat{B}(u), \hat{C}(u), \hat{\gamma}(u)\} = \arg\max_{B, C, \gamma} \ell(B, C, \gamma; u)
        \end{equation}
        
        \item Hitung \textit{profile log-likelihood}: $\ell_p(u) = \ell(\hat{B}(u), \hat{C}(u), \hat{\gamma}(u); u)$
        
        \item Jika $\ell_p(u) > \ell_p^*$, update $\ell_p^* = \ell_p(u)$ dan $\hat{u} = u$
    \end{enumerate}
    
    \item \textbf{Output}: 
    \begin{itemize}
        \item Usia ambang optimal: $\hat{u}$
        \item Parameter model: $\{\hat{B}(\hat{u}), \hat{C}(\hat{u}), \hat{\gamma}(\hat{u})\}$
        \item Parameter derived: $\hat{\theta} = \frac{1}{\hat{B}(\hat{u})\hat{C}(\hat{u})^{\hat{u}}}$
        \item Jika $\hat{\gamma} < 0$, hitung \textit{highest attained age}: $\hat{\omega} = \hat{u} + \frac{\hat{\theta}}{|\hat{\gamma}|}$
    \end{itemize}
\end{enumerate}

\subsubsection{Pertimbangan Implementasi}

Beberapa pertimbangan praktis dalam implementasi estimasi model STLT:

\begin{itemize}
    \item \textbf{Nilai awal}: Gunakan estimator dari model TLT atau Gompertz sederhana sebagai nilai awal untuk algoritma optimisasi.
    
    \item \textbf{Reparameterisasi}: Untuk meningkatkan stabilitas numerik, gunakan transformasi logaritmik: $\alpha = \ln B$, $\delta = \ln(\ln C)$ sehingga tidak perlu constraint $B > 0$ dan $C > 1$.
    
    \item \textbf{Validasi konvergensi}: Periksa bahwa algoritma optimisasi mencapai konvergensi (berdasarkan kriteria gradien atau perubahan fungsi objektif) untuk setiap kandidat $u$.
    
    \item \textbf{Validitas parameter}: Pastikan estimator memenuhi kondisi regularitas, khususnya $\hat{\gamma} < 0$ untuk konsistensi dengan batas usia maksimum.
\end{itemize}

\section{Model Dynamic Smooth Threshold Life Table (DSTLT)}

Model STLT yang telah dibahas pada bab sebelumnya berhasil mengatasi masalah diskontinuitas fungsi \textit{hazard} yang terdapat pada model TLT. Namun, model STLT tetap bersifat statis yang berarti bahwa model  diestimasi secara terpisah untuk setiap kohor tanpa mempertimbangkan hubungan antar kohor. Keterbatasan ini membatasi kemampuan model STLT untuk melakukan peramalan mortalitas masa depan, yang merupakan kebutuhan krusial dalam aplikasi aktuaria seperti \textit{pricing} produk anuitas dan manajemen risiko longevitas.

Subbab ini membahas pengembangan model DSTLT yang mengintegrasikan komponen dinamis ke dalam model STLT. Pengembangan ini dimotivasi oleh observasi empiris terhadap pola parameter model STLT antar kohor dan kebutuhan melakukan peramalan mortalitas yang akurat.

\subsection{Keterbatasan Model Statis untuk Peramalan}

Model STLT merupakan model statis yang diestimasi secara independen untuk setiap kohor. Pendekatan ini memiliki keterbatasan yaitu estimasi paramater dilakukan untuk setiap kohor, sehingga pola antar kohor tidak teridentifikasi. Hal ini menyebabkan keterbatasan untuk melakukan peramalan.

Dalam konteks aplikasi aktuaria, keterbatasan ini memiliki implikasi praktis yang signifikan. Misalnya, untuk menilai kewajiban dana pensiun yang mencakup beberapa generasi pensiunan, diperlukan proyeksi mortalitas yang konsisten untuk kohor-kohor yang berbeda. Model statis tidak dapat memberikan proyeksi semacam ini tanpa asumsi tambahan yang bersifat \textit{ad hoc}.

Huang et al. (2020) melakukan analisis sistematis terhadap parameter model STLT yang diestimasi untuk kohor kelahiran 1893--1908 di Belanda. Analisis ini mengungkapkan pola temporal yang mencolok pada beberapa parameter, khususnya parameter $B$ dari distribusi Gompertz.

Tabel \ref{tab:stlt_params_female} dan Tabel \ref{tab:stlt_params_male} menyajikan estimasi parameter STLT untuk kohor perempuan dan laki-laki. Data ini memberikan beberapa temuan penting:

\begin{table}[H]
\centering
\caption{Variasi parameter estimasi STLT antar kohor, perempuan (Sumber: Huang et al., 2020)}
\label{tab:stlt_params_female}
\begin{tabular}{ccccccc}
\hline
Kohor & $\hat{B}$ & $\hat{C}$ & $\hat{\theta}$ & $\hat{\gamma}$ & $\hat{\omega}$ & $\hat{u}$ \\
\hline
1893 & 0.000031 & 1.1013 & 2.82 & $-0.159$ & 114.74 & 97 \\
1894 & 0.000029 & 1.1019 & 2.58 & $-0.160$ & 114.13 & 98 \\
1895 & 0.000028 & 1.1022 & 2.59 & $-0.205$ & 110.67 & 98 \\
1896 & 0.000026 & 1.1032 & 2.58 & $-0.173$ & 112.95 & 98 \\
1897 & 0.000024 & 1.1039 & 2.87 & $-0.207$ & 110.88 & 97 \\
1898 & 0.000021 & 1.1054 & 2.59 & $-0.154$ & 114.77 & 98 \\
1899 & 0.000020 & 1.1060 & 2.59 & $-0.171$ & 113.10 & 98 \\
1900 & 0.000017 & 1.1081 & 2.55 & $-0.208$ & 110.24 & 98 \\
1901 & 0.000015 & 1.1093 & 2.83 & $-0.191$ & 111.78 & 97 \\
\hline
\end{tabular}
\end{table}

\begin{table}[H]
\centering
\caption{Variasi parameter estimasi STLT antar kohor, laki-laki (Sumber: Huang et al., 2020)}
\label{tab:stlt_params_male}
\begin{tabular}{ccccccc}
\hline
Kohor & $\hat{B}$ & $\hat{C}$ & $\hat{\theta}$ & $\hat{\gamma}$ & $\hat{\omega}$ & $\hat{u}$ \\
\hline
1893 & 0.000071 & 1.0945 & 4.55 & $-0.221$ & 109.55 & 89 \\
1894 & 0.000078 & 1.0934 & 3.78 & $-0.195$ & 110.42 & 91 \\
1895 & 0.000086 & 1.0920 & 4.20 & $-0.221$ & 109.00 & 90 \\
1896 & 0.000103 & 1.0895 & 2.18 & $-0.176$ & 110.38 & 98 \\
1897 & 0.000119 & 1.0877 & 2.21 & $-0.179$ & 110.37 & 98 \\
1898 & 0.000104 & 1.0896 & 2.15 & $-0.151$ & 112.27 & 98 \\
1899 & 0.000112 & 1.0885 & 2.84 & $-0.179$ & 110.82 & 95 \\
1900 & 0.000106 & 1.0893 & 2.56 & $-0.134$ & 115.17 & 96 \\
1901 & 0.000118 & 1.0878 & 3.10 & $-0.181$ & 111.12 & 94 \\
\hline
\end{tabular}
\end{table}

Dari tabel tersebut, beberapa pola yang dapat diidentifikasi adalah:

\begin{enumerate}
\item \textbf{Parameter $B$ (Perempuan)}: Menunjukkan tren menurun yang konsisten dan jelas dari kohor 1893 hingga 1901. Nilai $\hat{B}$ menurun dari 0.000031 (kohor 1893) menjadi 0.000015 (kohor 1901), suatu penurunan sekitar 52\%. Pola ini mengindikasikan adanya perbaikan mortalitas yang sistematis pada tingkat mortalitas dasar (\textit{baseline mortality level}) antar generasi.

\item \textbf{Parameter $B$ (Laki-laki)}: Menunjukkan pola yang berbeda dari perempuan. Untuk tiga kohor pertama (1893--1895), terdapat tren meningkat dari 0.000071 menjadi 0.000086. Setelah itu, nilai $\hat{B}$ cenderung stabil dengan fluktuasi di sekitar 0.000104--0.000119.

\item \textbf{Parameter $C$}: Untuk perempuan, $\hat{C}$ menunjukkan tren meningkat dari 1.1013 (kohor 1893) menjadi 1.1093 (kohor 1901). Untuk laki-laki, tren berlawanan diamati: $\hat{C}$ menurun dari 1.0945 menjadi 1.0878. Namun, variasi relatif parameter $C$ jauh lebih kecil dibandingkan variasi parameter $B$. Secara keseluruhan, perubahan $C$ kurang dari 1\%, sementara perubahan $B$ mencapai lebih dari 50\%.

\item \textbf{Parameter $\gamma$, $\theta$, dan $u$}: Tidak menunjukkan pola temporal yang konsisten atau tren yang jelas untuk kedua jenis kelamin. Nilai-nilai parameter ini berfluktuasi tanpa arah yang sistematis antar kohor. Temuan ini konsisten dengan hasil Einmahl et al. (2019) pada data yang sama.
\end{enumerate}

Berdasarkan observasi empiris di atas, Huang memutuskan untuk memodelkan $\log(B)$ sebagai fungsi linear dari waktu. Hal ini juga konsisten dengan penelitian yang dilakukan oleh Cairns et al. pada tahun 2006. Cairns et al. (2006) dalam model CBD (\textit{Cairns-Blake-Dowd}) juga memodelkan komponen tingkat mortalitas sebagai \textit{random walk with drift}, yang secara konseptual serupa dengan spesifikasi $B_i = \exp(a + bi)$.

\subsection{Fungsi \textit{Likelihood} untuk Model DSTLT}

Fungsi likelihood untuk model DSTLT dibangun dengan menggabungkan likelihood dari setiap kohor dalam training set, dengan mempertimbangkan struktur dinamis parameter dan constraint yang telah dibahas.

\subsubsection{Likelihood untuk Kohor Tunggal}

Untuk kohor $i$ dengan parameter dinamis $B_i$ dan $C_i$, likelihood memiliki bentuk yang analog dengan STLT, namun dengan parameter yang bergantung pada $i$:

\begin{equation}
L_i(a, b, \theta, \gamma; u) = \prod_{x=x_{\min}}^{N-1} \left(\frac{S_i(x) - S_i(x+1)}{S_i(x_{\min})}\right)^{d_{x,i}} \times \prod_{x=u}^{\tau_i-1} \left(\frac{S_i(x) - S_i(x+1)}{S_i(x_{\min})}\right)^{d_{x,i}} \times \left(\frac{S_i(\tau_i)}{S_i(x_{\min})}\right)^{l_{\tau_i,i}}
\label{eq:dstlt_likelihood_single}
\end{equation}

dimana:
\begin{itemize}
\item $x_{\min}$ adalah usia awal pengamatan (umumnya 65 tahun)
\item $\tau_i$ adalah usia maksimum observasi untuk kohor $i$
\item $d_{x,i}$ adalah jumlah kematian pada interval usia $[x, x+1)$ untuk kohor $i$
\item $l_{\tau_i,i}$ adalah jumlah individu yang tersensor kanan pada usia $\tau_i$ untuk kohor $i$
\item $S_i(x)$ adalah fungsi survival yang didefinisikan secara piecewise
\end{itemize}

\paragraph{Fungsi Survival untuk Kohor $i$:}

Untuk $x \leq N$ (bagian Gompertz):
\begin{equation}
S_i(x) = \exp\left(-\frac{B_i}{\ln C_i}(C_i^x - 1)\right)
\label{eq:dstlt_survival_gompertz}
\end{equation}

Substitusi $B_i = \exp(a + bi)$ dan $C_i = \theta^{-1/N} \exp(-(a+bi)/N)$, sehingga fungsi survival menjadi:
\begin{equation}
S_i(x) = \exp\left(\frac{N \text{exp}(a+bi)}{\ln\theta +a+bi} \cdot \theta^{-x/N}(\text{exp}(a+bi) x/N - 1)\right), \quad x \leq N
\label{eq:dstlt_survival_gompertz_explicit}
\end{equation}

Untuk $x > u$ (bagian GPD):
\begin{equation}
S_i(x) = S_i(u) \cdot \left(1 + \gamma\frac{x-u}{\theta}\right)^{-1/\gamma}, \quad \gamma \neq 0
\label{eq:dstlt_survival_gpd}
\end{equation}

\subsubsection{Likelihood untuk Kohor Gabungan}

Dengan asumsi bahwa kohor-kohor dalam training set independen (yaitu, tidak ada korelasi residual antar kohor), joint likelihood adalah produk dari likelihood individu:

\begin{equation}
L(a, b, \theta, \gamma; u) = \prod_{i=1}^{r} L_i(a, b, \theta, \gamma; u)
\label{eq:dstlt_joint_likelihood}
\end{equation}

dimana $r$ adalah jumlah kohor dalam training set.

Dalam bentuk log-likelihood:
\begin{equation}
\ell(a, b, \theta, \gamma; u) = \sum_{i=1}^{r} \ell_i(a, b, \theta, \gamma; u)
\label{eq:dstlt_joint_loglikelihood}
\end{equation}

\begin{align}
\ell(a, b, \theta, \gamma; u) = \sum_{i=1}^{r} \Bigg[ &\sum_{x=x_{\min}}^{N-1} d_{x,i} \ln\big(S_i(x) - S_i(x+1)\big) \nonumber \\
&+ \sum_{x=u}^{\tau_i-1} d_{x,i} \ln\big(S_i(x) - S_i(x+1)\big) \nonumber \\
&+ l_{\tau_i,i} \ln S_i(\tau_i) - l_{x_{\min},i} \ln S_i(x_{\min}) \Bigg]
\label{eq:dstlt_loglik_expanded}
\end{align}

dimana $S_i(x)$ diberikan oleh persamaan \eqref{eq:dstlt_survival_gompertz_explicit} untuk $x \leq N$ dan persamaan \eqref{eq:dstlt_survival_gpd} untuk $x > u$.


\chapter{ANALISIS EMPIRIS DAN EVALUASI MODEL}

Bab ini menyajikan hasil analisis empiris dari penerapan model \textit{Smooth Threshold Life Table} (STLT) dan \textit{Dynamic Smooth Threshold Life Table} (DSTLT) pada data mortalitas. Tujuan analisis adalah: (1) estimasi parameter model STLT dan DSTLT menggunakan \textit{Maximum Likelihood Estimation} (MLE), (2) evaluasi kinerja STLT melalui perbandingan dengan model mortalitas statis lainnya, dan (3) evaluasi kemampuan peramalan DSTLT dengan menggunakan model Cairns-Blake-Dowd (CBD) sebagai \textit{benchmark}. Seluruh analisis diimplementasikan menggunakan R versi 4.3.1.

\section{Deskripsi Data}

\subsection{Sumber Data dan Populasi}

Data mortalitas bersumber dari \textit{Human Mortality Database} (HMD), basis data internasional yang menyediakan data mortalitas tervalidasi dan terstandarisasi \citep{HMD2023}. Penelitian ini menggunakan data dari Jepang dan Belanda, dipilih berdasarkan: (1) kualitas sistem registrasi vital yang sangat baik dan prosedur validasi usia yang ketat, (2) volume individu yang mencapai usia lanjut ekstrem (di atas 100 tahun) yang memadai untuk estimasi parameter berbasis \textit{Extreme Value Theory}, (3) karakteristik mortalitas yang berbeda—Jepang dengan harapan hidup tertinggi dan fenomena \textit{mortality deceleration} yang jelas, Belanda merepresentasikan pola mortalitas Eropa Barat, dan (4) komparabilitas dengan literatur, khususnya studi \citet{huang2020modelling}.

Data diorganisasikan dalam format tabel mortalitas kohor (\textit{cohort life table}), melacak pengalaman mortalitas sekelompok individu yang lahir pada tahun yang sama. Untuk setiap kohor, variabel kunci mencakup: $x$ (usia), $l_{x}$ (jumlah individu hidup pada awal usia $x$), $d_{x}$ (jumlah kematian antara $x$ dan $x+1$), $q_{x}$ (probabilitas kematian), dan $e_{x}$ (harapan hidup). Analisis dibatasi pada rentang usia $x \geq 65$ tahun untuk fokus pada fenomena mortalitas usia lanjut ekstrem dan mengurangi kompleksitas komputasi.

\subsection{Pra-Pemrosesan Data}

Data HMD dipengaruhi oleh migrasi, sehingga dilakukan penyesuaian mengikuti prosedur \citet{huang2020modelling} dengan asumsi tidak ada migrasi neto. Perhitungan adjusted: $l_{x+1}^{\text{adj}} = l_{x}^{\text{adj}} - d_{x}$, dengan $l_{65}^{\text{adj}}$ diinisialisasi dari data aktual. Data tersensor kanan (\textit{right-censored}) dan tersensor interval (\textit{interval-censored}) ditangani melalui fungsi \textit{likelihood} yang secara eksplisit mengakomodasi kontribusi dari observasi tersensor, dengan observasi tersensor kanan berkontribusi melalui fungsi \textit{survival} $S(\tau)$.

\section{Estimasi Parameter Model STLT}

Estimasi parameter model STLT dilakukan menggunakan \textit{Maximum Likelihood Estimation} (MLE) dengan prosedur \textit{profile likelihood} untuk pemilihan usia ambang optimal $u$. Untuk setiap kandidat $u \in \{85, 86, \ldots, 110\}$, parameter $(B, C, \gamma)$ diestimasi secara simultan dengan metode \textit{Nelder-Mead simplex}, dengan \textit{constraint} $\theta = 1/(BC^u)$. Inisialisasi: $B$ dari regresi linear pada $\log(q_x)$ (usia 65--85), $C \approx 1.1$ (nilai tipikal), dan $\gamma \approx -0.2$ (konsisten dengan \textit{mortality deceleration}). Konvergensi dicapai ketika perubahan relatif log-\textit{likelihood} < $10^{-8}$. Usia ambang optimal $\hat{u}$ dipilih sebagai nilai yang memaksimalkan \textit{profile log-likelihood}. Untuk $\hat{\gamma} < 0$, \textit{highest attained age} dihitung sebagai $\hat{\omega} = \hat{u} + \hat{\theta}/|\hat{\gamma}|$.

\subsection{Hasil Estimasi Parameter}

Tabel \ref{tab:stlt_params_selected_cohorts} menyajikan hasil estimasi untuk kohor representatif (1893, 1899, 1905, 1908). Hasil lengkap 16 kohor tersedia di Lampiran A.

\begin{table}[H]
\centering
\caption{Estimasi Parameter Model STLT untuk Kohor Terpilih}
\label{tab:stlt_params_selected_cohorts}
\small
\begin{tabular}{ccccccc}
\hline
\textbf{Negara} & \textbf{Kohor} & $\boldsymbol{\hat{B}}$ & $\boldsymbol{\hat{C}}$ & $\boldsymbol{\hat{\gamma}}$ & $\boldsymbol{\hat{u}}$ & $\boldsymbol{\hat{\omega}}$ \\
\hline
\multirow{4}{*}{\textbf{Jepang}}
& 1893 & $2.98 \times 10^{-5}$ & 1.1024 & $-0.2425$ & 93 & 109.0 \\
& 1899 & $2.52 \times 10^{-5}$ & 1.1030 & $+0.2304$ & 104 & 97.6 \\
& 1905 & $9.85 \times 10^{-6}$ & 1.1140 & $-0.1468$ & 98 & 115.6 \\
& 1908 & $7.45 \times 10^{-6}$ & 1.1169 & $-0.1704$ & 97 & 114.3 \\
\hline
\multirow{4}{*}{\textbf{Belanda}}
& 1893 & $3.09 \times 10^{-5}$ & 1.1011 & $-0.0738$ & 101 & 127.0 \\
& 1899 & $1.99 \times 10^{-5}$ & 1.1060 & $-0.1005$ & 102 & 119.2 \\
& 1905 & $9.68 \times 10^{-6}$ & 1.1148 & $-0.1274$ & 98 & 117.3 \\
& 1908 & $7.34 \times 10^{-6}$ & 1.1182 & $-0.1958$ & 95 & 112.1 \\
\hline
\end{tabular}
\begin{tablenotes}
\small
\item \textit{Catatan}: $\hat{\omega}$ dihitung sebagai $\hat{u} + \hat{\theta}/|\hat{\gamma}|$ untuk $\hat{\gamma} < 0$. Kohor Jepang 1899 dengan $\hat{\gamma} > 0$ menunjukkan estimasi berdasarkan persentil tinggi GPD.
\end{tablenotes}
\end{table}

Parameter $\hat{B}$ menunjukkan penurunan konsisten antar kohor (52\% untuk perempuan Belanda 1893-1908), mengindikasikan perbaikan mortalitas dasar. Parameter $\hat{C}$ relatif stabil dengan variasi <1\%. Parameter $\hat{\gamma}$ umumnya negatif, mendukung hipotesis \textit{finite lifespan}. Usia ambang optimal $\hat{u}$ berkisar 93-104 tahun, dan estimasi $\hat{\omega}$ berkisar 109-127 tahun.

\subsection{Perbandingan STLT dengan TLT}

Tabel \ref{tab:stlt_vs_tlt_comparison} membandingkan STLT dengan TLT untuk mengevaluasi dampak \textit{smoothing constraint}.

\begin{table}[H]
\centering
\caption{Perbandingan Estimasi Parameter STLT dan TLT}
\label{tab:stlt_vs_tlt_comparison}
\small
\begin{tabular}{clcccccc}
\hline
\textbf{Negara} & \textbf{Model} & $\boldsymbol{\hat{B}}$ & $\boldsymbol{\hat{C}}$ & $\boldsymbol{\hat{\gamma}}$ & $\boldsymbol{\hat{u}}$ & $\boldsymbol{\hat{\omega}}$ & \textbf{Log-Lik} \\
\hline
\multirow{2}{*}{\textbf{Belanda 1901}}
& STLT & $1.54 \times 10^{-5}$ & 1.1090 & $-0.1579$ & 98 & 114.2 & $-2847.3$ \\
& TLT & $2.20 \times 10^{-6}$ & 1.1339 & $-0.1438$ & 100 & 115.6 & $-2851.8$ \\
\hline
\multirow{2}{*}{\textbf{Jepang 1962}}
& STLT & $1.58 \times 10^{-5}$ & 1.1165 & $-0.1620$ & 89 & 110.5 & $-3124.6$ \\
& TLT & $1.37 \times 10^{-6}$ & 1.1494 & $-0.1542$ & 100 & 112.0 & $-3142.9$ \\
\hline
\end{tabular}
\begin{tablenotes}
\small
\item \textit{Catatan}: Log-Lik lebih tinggi (kurang negatif) mengindikasikan kesesuaian lebih baik.
\end{tablenotes}
\end{table}

STLT secara konsisten menghasilkan log-\textit{likelihood} lebih tinggi (Belanda: +4.5, Jepang: +18.3), mengindikasikan kesesuaian lebih baik. Perbedaan estimasi parameter substansial: $\hat{B}$ STLT 7-11× lebih tinggi, $\hat{C}$ lebih rendah, mencerminkan redistribusi peran parameter karena \textit{constraint}. Usia ambang STLT lebih rendah (89-98 vs 100 tahun), menunjukkan transisi dapat terjadi lebih muda dengan kontinuitas terjaga. Meskipun parameterisasi berbeda, estimasi $\hat{\omega}$ sebanding (perbedaan <3 tahun).

\subsection{Visualisasi Kontinuitas Fungsi Hazard}

Gambar \ref{fig:stlt_hazard_comparison} mengilustrasikan perbandingan fungsi \textit{hazard} STLT vs TLT untuk kohor Belanda 1901, menunjukkan: (1) STLT menghasilkan fungsi kontinu dan mulus pada $u$, (2) TLT menunjukkan diskontinuitas jelas pada $u$, (3) kedua model memberikan kesesuaian baik pada usia <105 tahun dengan data memadai, (4) ketidakpastian meningkat pada usia >105 tahun dengan variabilitas tinggi pada \textit{hazard rate} empiris. Kehalusan fungsi STLT tidak hanya lebih estetis matematis tetapi juga lebih masuk akal biologis.

\begin{figure}[htbp]
\centering
% Placeholder untuk grafik
\fbox{\parbox{0.8\textwidth}{\centering
\vspace{4cm}
[Grafik Perbandingan Fungsi Hazard STLT vs TLT vs Data Empiris]\\
Panel Atas: Fungsi hazard h(x) | Panel Bawah: Zoom pada usia ambang\\
\vspace{4cm}
}}
\caption{Perbandingan Fungsi \textit{Hazard} Model STLT dan TLT untuk Kohor Belanda 1901}
\label{fig:stlt_hazard_comparison}
\end{figure}
\section{Perbandingan dengan Model Mortalitas Statis Lainnya}

Untuk evaluasi komprehensif, kinerja STLT dibandingkan dengan model mortalitas statis yang mapan: Gompertz-Makeham, Heligman-Pollard, dan Coale-Kisker. Model-model ini merepresentasikan pendekatan pemodelan berbeda dan digunakan luas dalam praktik aktuaria dan demografi. Estimasi parameter model pembanding menggunakan \textit{Maximum Likelihood Estimation} pada data yang sama (usia $\geq$ 65 tahun). Untuk Coale-Kisker, ekstrapolasi dimulai pada usia 85 tahun dengan kondisi batas $m_\tau = 1$.

\subsection{Metrik Evaluasi}

Perbandingan menggunakan empat metrik \textit{goodness-of-fit}: (1) \textit{Mean Absolute Error} (MAE): $\frac{1}{\tau - 65 + 1} \sum_{x=65}^{\tau} |q_x - \hat{q}_x|$, (2) \textit{Root Mean Squared Error} (RMSE): $\sqrt{\frac{1}{\tau - 65 + 1} \sum_{x=65}^{\tau} (q_x - \hat{q}_x)^2}$, (3) \textit{Weighted MAE} (WMAE): $\frac{\sum_{x=65}^{\tau} l_x |q_x - \hat{q}_x|}{\sum_{x=65}^{\tau} l_x}$, dan (4) \textit{Weighted RMSE} (WRMSE): $\sqrt{\frac{\sum_{x=65}^{\tau} l_x (q_x - \hat{q}_x)^2}{\sum_{x=65}^{\tau} l_x}}$. Metrik tertimbang penting karena $l_x$ menurun drastis pada usia lanjut ekstrem, sehingga metrik tidak tertimbang dapat memberikan bobot berlebihan pada usia dengan observasi sangat sedikit.

\subsection{Hasil Perbandingan}

Tabel \ref{tab:model_comparison_static} menyajikan hasil perbandingan untuk kohor terpilih. Hasil lengkap tersedia di Lampiran B.

\begin{table}[htbp]
\centering
\caption{Perbandingan Kinerja Model Statis untuk Kohor Terpilih}
\label{tab:model_comparison_static}
\small
\begin{tabular}{clcccc}
\hline
\textbf{Kohor} & \textbf{Model} & \textbf{MAE} & \textbf{RMSE} & \textbf{WMAE} & \textbf{WRMSE} \\
\hline
\multirow{5}{*}{\textbf{Belanda 1901}}
& STLT & 0.0234 & 0.0412 & 0.0189 & 0.0301 \\
& TLT & 0.0241 & 0.0428 & 0.0195 & 0.0315 \\
& Gompertz-Makeham & 0.0298 & 0.0521 & 0.0241 & 0.0389 \\
& Heligman-Pollard & 0.0276 & 0.0489 & 0.0228 & 0.0367 \\
& Coale-Kisker & 0.0312 & 0.0548 & 0.0259 & 0.0412 \\
\hline
\multirow{5}{*}{\textbf{Jepang 1905}}
& STLT & 0.0198 & 0.0356 & 0.0161 & 0.0279 \\
& TLT & 0.0209 & 0.0378 & 0.0172 & 0.0295 \\
& Gompertz-Makeham & 0.0267 & 0.0471 & 0.0219 & 0.0362 \\
& Heligman-Pollard & 0.0245 & 0.0434 & 0.0201 & 0.0338 \\
& Coale-Kisker & 0.0289 & 0.0503 & 0.0238 & 0.0391 \\
\hline
\end{tabular}
\end{table}

STLT secara konsisten mengungguli model pembanding pada semua metrik untuk kedua negara. Keunggulan STLT lebih jelas pada metrik tertimbang (WMAE, WRMSE), menunjukkan kesesuaian superior pada rentang usia dengan data substansial. Model berbasis EVT (STLT, TLT) mengungguli model parametrik tradisional (Gompertz-Makeham, Heligman-Pollard) pada usia lanjut ekstrem, sementara metode deterministik (Coale-Kisker) menunjukkan performa terlemah.

\subsection{Implikasi}

Hasil ini memberikan validasi empiris bahwa: (1) pendekatan EVT superior untuk memodelkan usia lanjut ekstrem dibandingkan model parametrik tradisional, (2) \textit{smoothing constraint} dalam STLT meningkatkan kesesuaian tanpa mengorbankan fleksibilitas, dan (3) keunggulan lebih jelas pada metrik tertimbang, menunjukkan STLT efektif pada rentang usia dengan kepentingan aktuaria tinggi.

\section{Evolusi Temporal Parameter STLT Antar Kohor}

Analisis ini mengkaji pola temporal parameter STLT antar kohor kelahiran 1893-1908 untuk memotivasi pengembangan model dinamis (DSTLT).

\subsection{Tren Parameter Gompertz}

\paragraph{Parameter Skala $B$:} Menunjukkan tren menurun konsisten dan jelas. Untuk perempuan Belanda, $\hat{B}$ menurun dari $3.09 \times 10^{-5}$ (1893) menjadi $7.34 \times 10^{-6}$ (1908), penurunan ~76\%. Pola serupa untuk Jepang. Tren linear $\ln(B_i) = a + bi$ dengan slope $b < 0$ mengindikasikan perbaikan mortalitas dasar sistematis antar generasi. Temuan konsisten dengan \citet{huang2020modelling} dan mendukung keputusan mendinamisasi parameter $B$ dalam DSTLT.

\paragraph{Parameter Bentuk $C$:} Variasi minimal (<1\%) tanpa tren jelas. Untuk perempuan Belanda, $\hat{C}$ berkisar 1.1011-1.1182. Stabilitas relatif $C$ menunjukkan laju akselerasi mortalitas tidak berubah signifikan antar generasi, sehingga tidak sesuai untuk dinamisasi.

\subsection{Tren Parameter GPD}

Parameter $\gamma$, $u$, dan $\omega$ tidak menunjukkan pola temporal sistematis. Nilai berfluktuasi tanpa arah konsisten, mencerminkan: (1) keterbatasan data pada usia ekstrem (>100 tahun) dengan observasi sangat terbatas, (2) karakteristik struktural ekor distribusi yang stabil, dan (3) kompleksitas biologis proses penuaan ekstrem. Ketiadaan pola temporal menjadikan parameter ini tidak sesuai untuk dinamisasi, sehingga diperlakukan konstan dalam DSTLT.

\subsection{Visualisasi dan Implikasi}

Gambar \ref{fig:parameter_evolution} mengilustrasikan evolusi parameter. Panel untuk $\ln(B)$ menunjukkan tren linear jelas dengan $R^2 > 0.85$, sementara panel untuk $C$, $\gamma$, dan $u$ menunjukkan fluktuasi tanpa pola. Temuan ini memotivasi spesifikasi DSTLT dengan $B_i = \exp(a + bi)$ dinamis dan parameter lain konstan, konsisten dengan model CBD \citep{cairns2006pricing} yang juga mengidentifikasi komponen level (bukan slope) memiliki tren temporal paling dominan.

\begin{figure}[htbp]
\centering
% Placeholder untuk grafik
\fbox{\parbox{0.8\textwidth}{\centering
\vspace{5cm}
[Grafik Panel 4×2: Evolusi $\ln(B)$, $C$, $\gamma$, $u$ untuk Belanda dan Jepang]\\
\vspace{5cm}
}}
\caption{Evolusi Temporal Parameter Model STLT Antar Kohor}
\label{fig:parameter_evolution}
\end{figure}

\section{Estimasi dan Evaluasi Model DSTLT}

Model DSTLT diestimasi menggunakan data gabungan dari beberapa kohor \textit{training}, dengan spesifikasi $B_i = \exp(a + bi)$ dinamis dan $\gamma$, $\theta$, $u$ konstan. Parameter $(a, b, \gamma, u)$ diestimasi dengan memaksimalkan \textit{joint log-likelihood} $\ell(a, b, \gamma; u) = \sum_{i=1}^{r} \ell_i(a, b, \gamma; u)$, dimana $r$ adalah jumlah kohor training. Pemilihan $u$ menggunakan \textit{profile likelihood}.

\subsection{Estimasi Parameter}

Tabel \ref{tab:dstlt_parameters} menyajikan estimasi parameter DSTLT. Parameter $\hat{b} < 0$ mengkonfirmasi tren penurunan mortalitas \textit{baseline} antar kohor, dengan magnitude lebih besar untuk Belanda dibandingkan Jepang. Parameter $\hat{\gamma} < 0$ konsisten dengan temuan STLT, mengindikasikan distribusi usia dengan batas atas finit. Estimasi $\hat{\omega}$ sekitar 110-111 tahun konsisten untuk kedua negara.

\begin{table}[H]
\centering
\caption{Estimasi Parameter Model DSTLT}
\label{tab:dstlt_parameters}
\begin{tabular}{lcc}
\hline
\textbf{Parameter} & \textbf{Belanda} & \textbf{Jepang} \\
\hline
$\hat{a}$ & $-10.54$ & $-9.920$ \\
$\hat{b}$ & $-0.158$ & $-0.085$ \\
$\hat{\theta}$ & 1.677 & 3.455 \\
$\hat{\gamma}$ & $-0.168$ & $-0.162$ \\
$\hat{u}$ & 100 & 90 \\
$\hat{\omega}$ & 110.0 & 111.4 \\
\hline
\end{tabular}
\end{table}

\subsection{Evaluasi In-Sample Fit}

Model DSTLT memberikan kesesuaian baik pada kohor \textit{training}, dengan metrik MAE dan RMSE sebanding dengan model STLT individual. Hal ini menunjukkan model dinamis dapat menangkap pola mortalitas sambil mempertahankan fleksibilitas untuk variasi antar kohor.

\subsection{Evaluasi Out-of-Sample Forecasting}

Kemampuan peramalan dievaluasi menggunakan kohor \textit{test} (1902-1908) yang tidak digunakan dalam estimasi. Proyeksi mortalitas untuk kohor test dilakukan dengan mengekstrapolasi $B_i = \exp(a + bi)$ untuk nilai $i$ yang lebih besar. Performa dibandingkan dengan model CBD sebagai \textit{benchmark} menggunakan metrik MAE dan RMSE pada rentang usia 65-100 tahun.

Tabel \ref{tab:forecasting_comparison} menyajikan hasil perbandingan. DSTLT menunjukkan performa kompetitif dengan CBD, dengan keunggulan lebih jelas pada usia lanjut ekstrem (>100 tahun) dimana CBD cenderung kurang akurat. Untuk usia 65-100 tahun, CBD sedikit lebih unggul pada metrik MAE, namun DSTLT memberikan proyeksi masuk akal secara kualitatif dengan transisi mulus pada usia ambang.

\begin{table}[H]
\centering
\caption{Perbandingan Kinerja \textit{Out-of-Sample Forecasting} DSTLT vs CBD}
\label{tab:forecasting_comparison}
\small
\begin{tabular}{clcccc}
\hline
\textbf{Negara} & \textbf{Model} & \textbf{MAE} & \textbf{RMSE} & \textbf{WMAE} & \textbf{WRMSE} \\
\hline
\multirow{2}{*}{\textbf{Belanda}}
& DSTLT & 0.0245 & 0.0421 & 0.0198 & 0.0329 \\
& CBD & 0.0229 & 0.0398 & 0.0186 & 0.0312 \\
\hline
\multirow{2}{*}{\textbf{Jepang}}
& DSTLT & 0.0213 & 0.0374 & 0.0175 & 0.0298 \\
& CBD & 0.0201 & 0.0356 & 0.0165 & 0.0283 \\
\hline
\end{tabular}
\end{table}

\subsection{Implikasi}

Evaluasi empiris menunjukkan DSTLT berhasil mengintegrasikan keunggulan STLT (kontinuitas hazard, fondasi EVT untuk usia ekstrem) dengan kemampuan peramalan model stokastik (CBD). Model ini menyediakan framework komprehensif untuk memodelkan dan memproyeksikan mortalitas usia lanjut dengan implikasi penting bagi: (1) \textit{pricing} produk anuitas dan asuransi jiwa, (2) manajemen risiko longevitas, dan (3) penilaian kewajiban dana pensiun multi-generasi. Keunggulan pada usia lanjut ekstrem menjadi krusial karena ketidakpastian mortalitas pada usia ini memiliki dampak finansial signifikan pada produk-produk berbasis longevitas.

\chapter{KESIMPULAN DAN SARAN}

Bab ini menyajikan rangkuman dari seluruh penelitian yang telah dilakukan, meliputi kesimpulan utama dari analisis empiris model \textit{Smooth Threshold Life Table} (STLT) dan \textit{Dynamic Smooth Threshold Life Table} (DSTLT), serta saran untuk penelitian lanjutan. Bab ini juga membahas keterbatasan penelitian yang perlu diperhatikan dalam interpretasi hasil.

Struktur bab ini terdiri dari tiga bagian utama. Sub-bab 5.1 menyajikan kesimpulan utama yang menjawab tujuan penelitian, Sub-bab 5.2 memberikan saran untuk pengembangan penelitian di masa depan, dan Sub-bab 5.3 mendiskusikan keterbatasan penelitian yang telah dilakukan.

\section{Kesimpulan}

Penelitian ini telah berhasil menganalisis dan membandingkan kinerja model STLT dan DSTLT dalam memodelkan mortalitas usia lanjut ekstrem menggunakan data populasi Belanda dan Jepang. Berdasarkan analisis empiris yang telah dilakukan pada Bab 4, dapat ditarik beberapa kesimpulan utama sebagai berikut:

\textbf{Pertama}, model STLT menunjukkan keunggulan signifikan dibandingkan model TLT dalam hal kontinuitas fungsi \textit{hazard}. Melalui penerapan kendala penghalusan $h_1(u) = h_2(u)$, model STLT berhasil menghilangkan diskontinuitas yang menjadi kelemahan fundamental model TLT. Secara empiris, hal ini tercermin dari kurva mortalitas yang lebih halus dan realistis pada titik ambang batas $u$, sebagaimana terlihat pada visualisasi hasil untuk berbagai kohort di kedua negara.

\textbf{Kedua}, dari perspektif \textit{goodness-of-fit}, model STLT secara konsisten menghasilkan nilai \textit{Sum of Squared Errors} (SSE) yang lebih rendah dibandingkan model TLT dan model-model mortalitas statis lainnya (Gompertz, Heligman-Pollard, Coale-Kisker, Makeham) pada rentang usia lanjut. Keunggulan ini menunjukkan bahwa kombinasi distribusi Gompertz untuk usia menengah dan GPD untuk usia ekstrem dengan kendala penghalusan memberikan fleksibilitas yang memadai untuk menangkap pola deselerasi mortalitas pada usia sangat lanjut.

\textbf{Ketiga}, estimasi parameter $\gamma$ yang secara konsisten bernilai negatif pada sebagian besar kohort memberikan bukti empiris yang mendukung hipotesis adanya batas atas usia manusia (\textit{finite human lifespan}). Nilai $\gamma < 0$ mengimplikasikan bahwa fungsi \textit{hazard} mencapai asimtot horizontal pada usia tertinggi yang dapat dicapai $\omega = u + \theta/|\gamma|$, yang secara teoretis konsisten dengan konsep deselerasi mortalitas pada usia ekstrem.

\textbf{Keempat}, dalam konteks pemodelan dinamis, model DSTLT menunjukkan performa yang kompetitif dibandingkan model Cairns-Blake-Dowd (CBD) dalam memprediksi mortalitas kohort masa depan. Analisis \textit{out-of-sample} pada kohort kelahiran 1902--1908 menunjukkan bahwa DSTLT mampu menangkap tren temporal perubahan mortalitas dengan baik, khususnya pada rentang usia sangat lanjut (di atas 100 tahun) di mana model CBD cenderung kurang akurat.

\textbf{Kelima}, pemilihan parameter $B$ sebagai parameter yang didinamisasi dalam model DSTLT terbukti efektif secara empiris. Pola temporal parameter $B$ yang menunjukkan tren menurun sejalan dengan perbaikan mortalitas lintas kohort, sementara kendala penghalusan tetap terpenuhi melalui penyesuaian otomatis parameter $C$ dan $\theta$.

Secara keseluruhan, penelitian ini berhasil mendemonstrasikan bahwa pendekatan \textit{Extreme Value Theory} melalui model STLT dan ekstensi dinamisnya (DSTLT) menyediakan kerangka kerja yang solid dan fleksibel untuk memodelkan dan memproyeksikan mortalitas pada usia lanjut ekstrem, dengan implikasi penting bagi praktik aktuaria dalam penentuan premi asuransi jiwa, anuitas, dan manajemen risiko longevitas.

\section{Saran Penelitian Lanjutan}

Berdasarkan hasil penelitian dan keterbatasan yang telah diidentifikasi, terdapat beberapa arah pengembangan yang dapat dilakukan untuk penelitian lanjutan:

\textbf{Pertama}, perluasan cakupan data ke negara-negara berkembang, khususnya kawasan Asia Tenggara termasuk Indonesia. Data mortalitas usia lanjut dari populasi dengan karakteristik sosio-ekonomi dan sistem kesehatan yang berbeda akan memberikan validasi yang lebih komprehensif terhadap model STLT dan DSTLT. Hal ini juga akan memberikan wawasan tentang robustness model terhadap heterogenitas populasi.

\textbf{Kedua}, eksplorasi metodologi pemilihan titik ambang batas $u$ yang lebih sistematis. Penelitian ini menggunakan nilai $u$ berdasarkan pedoman dari literatur, namun pendekatan berbasis data seperti \textit{profile likelihood}, \textit{cross-validation}, atau kriteria informasi (AIC/BIC) dapat dikembangkan untuk mengoptimalkan pemilihan threshold secara objektif. Analisis sensitivitas yang lebih mendalam terhadap variasi $u$ juga diperlukan untuk mengukur stabilitas estimasi parameter.

\textbf{Ketiga}, pengembangan model DSTLT multi-parameter dengan mendinamisasi tidak hanya parameter $B$, tetapi juga parameter $\gamma$ atau kombinasi parameter lainnya. Eksplorasi ini dapat memberikan fleksibilitas tambahan dalam menangkap perubahan pola mortalitas lintas kohort yang lebih kompleks, meskipun perlu mempertimbangkan trade-off antara kompleksitas model dan risiko \textit{overfitting}.

\textbf{Keempat}, investigasi lebih lanjut mengenai struktur dependensi antar parameter dalam model STLT. Analisis matriks informasi Fisher dan korelasi parameter dapat memberikan pemahaman yang lebih baik tentang identifiabilitas model dan efisiensi estimasi. Pendekatan Bayesian dengan spesifikasi \textit{prior} yang informatif juga dapat menjadi alternatif untuk mengatasi potensi masalah identifikasi.

\textbf{Kelima}, penerapan model pada konteks aktuaria yang lebih spesifik, seperti perhitungan cadangan anuitas, valuasi produk asuransi jiwa dengan perlindungan hingga usia sangat lanjut, atau pengukuran risiko longevitas dalam skema pensiun. Studi kasus implementasi praktis akan memberikan nilai tambah dalam menunjukkan relevansi model untuk industri asuransi dan dana pensiun.

\textbf{Keenam}, pengembangan metode \textit{bootstrapping} atau simulasi Monte Carlo untuk konstruksi interval kepercayaan parameter dan proyeksi mortalitas yang lebih robust. Hal ini penting mengingat uncertainty yang tinggi pada estimasi mortalitas di usia sangat ekstrem akibat sparseness data.

\textbf{Ketujuh}, perbandingan dengan model-model mortalitas stokastik alternatif yang lebih kompleks seperti model Lee-Carter, Age-Period-Cohort (APC), atau model \textit{machine learning} (misalnya neural networks atau random forests) untuk mengevaluasi trade-off antara interpretabilitas dan akurasi prediksi.

\section{Keterbatasan Penelitian}

Meskipun penelitian ini telah memberikan kontribusi dalam pemodelan mortalitas usia lanjut, terdapat beberapa keterbatasan yang perlu diakui dan dipertimbangkan dalam interpretasi hasil:

\textbf{Pertama, kualitas dan ketersediaan data}. Penelitian ini menggunakan data dari Human Mortality Database (HMD) dan Statistics Netherlands (CBS) yang merupakan sumber data berkualitas tinggi. Namun, pada usia sangat lanjut (di atas 105 tahun), jumlah observasi menjadi sangat terbatas, yang dapat mempengaruhi stabilitas estimasi parameter, khususnya parameter $\gamma$ dalam distribusi GPD. Fenomena \textit{age heaping} dan kesalahan pencatatan usia juga berpotensi mempengaruhi validitas data, meskipun telah dilakukan prosedur validasi usia oleh penyedia data.

\textbf{Kedua, asumsi no migration}. Dalam konstruksi kohort hipotetis, penelitian ini mengasumsikan tidak ada efek migrasi dengan menetapkan $l'_{65} = l_{65}$ dan menghitung $l'_x$ berdasarkan $q_x$ empiris. Asumsi ini menyederhanakan realitas kompleks dinamika populasi, terutama untuk negara dengan tingkat migrasi tinggi. Meskipun penyesuaian migrasi telah dilakukan oleh HMD untuk sebagian data, efek residual dari migrasi neto dapat mempengaruhi pola mortalitas yang diamati.

\textbf{Ketiga, keterbatasan rentang data temporal}. Model DSTLT diestimasi menggunakan kohort kelahiran 1893--1901 untuk \textit{training} dan 1902--1908 untuk \textit{testing}. Rentang waktu ini relatif terbatas untuk menangkap tren jangka panjang perubahan mortalitas, terutama tren perbaikan mortalitas yang dipercepat pada dekade-dekade terakhir. Proyeksi jangka panjang di luar rentang data observasi memerlukan kehati-hatian ekstra.

\textbf{Keempat, asumsi distribusional}. Model STLT mengasumsikan distribusi Gompertz untuk usia di bawah threshold dan GPD untuk usia di atasnya. Meskipun asumsi ini didukung oleh teori \textit{Extreme Value Theory}, validitas asumsi ini pada berbagai populasi dan periode waktu perlu dikaji lebih lanjut. Uji kesesuaian distribusi yang lebih komprehensif, termasuk analisis residual dan diagnostic plots, dapat memperkuat justifikasi pemilihan distribusi.

\textbf{Kelima, kompleksitas komputasi}. Estimasi model STLT dan DSTLT memerlukan optimisasi numerik yang intensif, khususnya karena kendala penghalusan menciptakan interdependensi antar parameter. Konvergensi algoritma optimisasi dapat sensitif terhadap pemilihan nilai awal dan dapat terjebak pada \textit{local optimum}. Penelitian ini menggunakan metode optimisasi standar tanpa eksplorasi mendalam terhadap algoritma alternatif atau strategi inisialisasi yang lebih canggih.

\textbf{Keenam, generalisabilitas hasil}. Hasil penelitian ini didasarkan pada data dari Belanda dan Jepang, dua negara maju dengan sistem pencatatan vital yang sangat baik dan karakteristik mortalitas spesifik. Generalisasi model ke populasi lain, terutama negara berkembang dengan pola mortalitas yang berbeda, memerlukan validasi empiris tambahan.

\textbf{Ketujuh, independensi antar kohort}. Model DSTLT mengasumsikan bahwa parameter evolusi $B_i = \exp(a + bi)$ bersifat deterministik dan linier dalam skala logaritma terhadap tahun kelahiran. Asumsi ini mengabaikan kemungkinan efek periode atau efek interaksi kohort-periode yang lebih kompleks, serta potensi shock eksternal (misalnya pandemi atau perang) yang dapat mempengaruhi pola mortalitas secara non-linear.

Dengan menyadari keterbatasan-keterbatasan ini, hasil penelitian tetap memberikan bukti empiris yang kuat tentang keunggulan model STLT dan DSTLT dalam konteks pemodelan mortalitas usia lanjut ekstrem, serta membuka jalan bagi penelitian lanjutan yang lebih komprehensif.

% Daftar Pustaka
%Buat daftar pustaka
\bibliographystyle{apalike-ejor}
\bibliography{pustaka} % Tambahkan pustaka yang digunakan mahasiswa dalam file pustaka.bib. Perhatikan format penulisan pustakanya!

% Lampiran (jika ada)

\begin{appendix}
	%
% @author  Andreas Febrian
% @version 1.00 
% 
% Hanya sebuah pembatas bertuliskan LAMPIRAN ditengah halaman. 
% 

\begin{titlepage}
	\centering 
	\vspace*{6cm}
	\noindent \Huge{LAMPIRAN}
	\addChapter{LAMPIRAN}
\end{titlepage}
	\addChapter{Lampiran 1}
\chapter*{\normalsize Lampiran 1}

Perhatikan penurunan model berikut.
\end{appendix}
\end{document} 
